\documentclass[11pt]{article}
\usepackage[margin=1in,headheight=15pt]{geometry}
\usepackage{fontspec}
\IfFileExists{cmunrm.otf}{\setmainfont{cmunrm.otf}[BoldFont=cmunbx.otf,ItalicFont=cmunti.otf,BoldItalicFont=cmunbi.otf]}{\setmainfont{lmroman10-regular.otf}[BoldFont=lmroman10-bold.otf,ItalicFont=lmroman10-italic.otf,BoldItalicFont=lmroman10-bolditalic.otf]}
\IfFileExists{cmuntt.otf}{\setmonofont{cmuntt.otf}[Scale=0.86,BoldFont=cmuntb.otf,ItalicFont=cmunit.otf,BoldItalicFont=cmuntx.otf]}{\setmonofont{lmmono10-regular.otf}[Scale=0.86]}
\usepackage{amsmath,amssymb,amsthm,mathtools}
\usepackage{booktabs,longtable,array,enumitem,microtype}
\usepackage{xcolor,hyperref,fancyhdr,listings}
\definecolor{linkblue}{RGB}{30,65,110}
\hypersetup{colorlinks=true,linkcolor=linkblue,citecolor=linkblue,urlcolor=linkblue,
 pdftitle={Combinatorial Transseries and Their Inverses: q-Products, Finite Fields, Theta Sectors, and Arithmetic Sheets},
 pdfauthor={Research extension prepared for Vladimir Reshetnikov}}
\pagestyle{fancy}\fancyhf{}
\fancyhead[L]{\small Combinatorial transseries and their inverses}
\fancyhead[R]{\small\thepage}
\renewcommand{\headrulewidth}{0.3pt}
\setlength{\parindent}{1.25em}
\setlength{\parskip}{3pt}
\setlength{\emergencystretch}{2em}
\allowdisplaybreaks[2]
\numberwithin{equation}{section}
\newtheorem{theorem}{Theorem}[section]
\newtheorem{proposition}[theorem]{Proposition}
\newtheorem{lemma}[theorem]{Lemma}
\newtheorem{corollary}[theorem]{Corollary}
\theoremstyle{definition}
\newtheorem{definition}[theorem]{Definition}
\newtheorem{example}[theorem]{Example}
\newtheorem{remark}[theorem]{Remark}
\newcommand{\N}{\mathbb N}
\newcommand{\Z}{\mathbb Z}
\newcommand{\R}{\mathbb R}
\newcommand{\C}{\mathbb C}
\newcommand{\F}{\mathbb F}
\newcommand{\GL}{\operatorname{GL}}
\newcommand{\Sp}{\operatorname{Sp}}
\newcommand{\Li}{\operatorname{Li}}
\newcommand{\rad}{\operatorname{rad}}
\newcommand{\coef}[2]{[#1^{#2}]}
\newcommand{\qbin}[2]{\genfrac{[}{]}{0pt}{}{#1}{#2}_{q}}
\newcommand{\qbinQ}[2]{\genfrac{[}{]}{0pt}{}{#1}{#2}_{Q}}
\newcommand{\pp}[1]{(Q;Q)_{#1}}
\newcommand{\Pinf}{\mathcal P_Q}
\newcommand{\A}{\mathcal A_Q}
\newcommand{\ee}{\mathrm e}
\newcommand{\dd}{\mathrm d}
\newcommand{\eps}{\varepsilon}
\newcommand{\T}{\Theta}
\newcommand{\bigO}{\mathcal O}
\newcommand{\one}{\mathbf 1}
\lstset{basicstyle=\ttfamily\footnotesize,breaklines=true,columns=fullflexible,
 frame=single,framerule=0.25pt,keepspaces=true,showstringspaces=false}

\begin{document}
\begin{titlepage}
\centering
\vspace*{1.1cm}
{\Large An extension of \emph{Transseries and Inversion}\par}
\vspace{1.1cm}
{\huge\bfseries Combinatorial Transseries\\[0.2em] and Their Inverses\par}
\vspace{0.65cm}
{\LARGE $q$-Products, Finite Fields, Theta Sectors,\\[0.2em] and Arithmetic Sheets\par}
\vspace{1cm}
{\large Prepared for Vladimir Reshetnikov\par}
\vspace{0.3cm}
{\large 4 September 2026\par}
\vspace{1cm}
\begin{minipage}{0.91\textwidth}
\begin{abstract}
This article develops a complementary collection of combinatorial examples in which the exponentially small structure can often be computed exactly, rather than inferred from a divergent power expansion. The main families are general linear, symplectic, and unitary group orders; complete and partial flag counts; fixed-rank and proportionally scaled Gaussian binomial coefficients; binomial-version $q$-Catalan and $q$-Fuss--Catalan numbers; and ordinary and generalized Galois numbers. For the latter, an entire generating function for all exponential corrections is obtained from theta series of the $A_{r-1}$ root lattice.

A common inversion theorem turns an analytic exponential tail attached to a quadratic or linear phase into a \emph{convergent} inverse transseries. Every coefficient is given by a finite, nonrecursive sum. A second inversion construction treats primitive necklaces and irreducible polynomials on fixed-radical subsequences, where the leading inverse is expressed by $W_{-1}$. Uniform integer-threshold bounds remain available even when no single smooth interpolation is selected. Finally, the limit $q\to1$ is analyzed in a simultaneous large-index scaling; the dominant phase changes to a dilogarithmic function, and a separate modular exponential scale appears.

All main identities and inversion assertions are proved. The accompanying Python program checks exact enumeration, the theta expansion, inverse coefficients through weight eight, and arithmetic threshold decisions, and reproduces the numerical tables.
\end{abstract}
\end{minipage}
\vfill
{\small Standalone companion article\;\textbullet\; explicit interpolation conventions\;\textbullet\; reproducible computations}
\end{titlepage}
\tableofcontents
\clearpage

\section{Scope, source relationship, and the families treated}
\label{sec:scope}

The linked \emph{Transseries and Inversion} article has become a substantial consolidated treatment of combinatorial threshold problems \cite{base}. Its examples include factorial-type functions, Bell and Fubini numbers, partitions, and rooted trees, among many others. The present extension follows the same central distinction between a leading inverse center, a full continuous inverse, and an integer threshold. Rather than rederive those examples, we organize a different collection around exact $q$-products, finite-field enumeration, root-lattice theta functions, and arithmetic divisor sums.

There are three particularly useful features of these families. First, their normalized tails are frequently analytic in $q^{-x}$, so that their transseries are genuinely convergent for large $x$. Second, inversion produces coefficients with a very explicit finite combinatorial structure. Third, two different kinds of discrete modulation occur: a finite number of residue-class theta constants, and an unbounded collection of divisibility indicators. These should not be conflated.

Throughout the fixed-$q$ sections,
\begin{equation}\label{eq:qnotation}
 q>1,\qquad Q=q^{-1}\in(0,1),\qquad h=\log q=-\log Q.
\end{equation}
A finite-field interpretation requires $q$ to be a prime power. The analytic formulas themselves are meaningful for every real $q>1$. A separate subsection treats $q$-factorials with base in $(0,1)$.

\begin{longtable}{@{}>{\raggedright\arraybackslash}p{0.25\textwidth}>{\raggedright\arraybackslash}p{0.32\textwidth}>{\raggedright\arraybackslash}p{0.34\textwidth}@{}}
\caption{The principal examples and their inverse scales. Constants and residue labels are retained in the exact centers used later.}\label{tab:map}\\
\toprule
Family & Leading size & Structure of the inverse correction\\
\midrule\endfirsthead
\toprule Family & Leading size & Structure of the inverse correction\\\midrule\endhead
$\lvert\GL(n,q)\rvert$ & $\Pinf q^{n^2}$ & Powers of $q^{-s}$, with $s\asymp\sqrt{\log Y}$\\
Complete flags, $[n]_q!$ & $\Pinf q^{n(n+1)/2}(q-1)^{-n}$ & Powers of $q^{-s}$ around an exact quadratic center\\
Fixed-rank $\qbin{n}{k}$ & $q^{k(n-k)}/\pp{k}$ & A convergent Puiseux series in $Y^{-1/k}$\\
Scaled Gaussian coefficients and partial flags & $Cq^{a_0n^2}$ & An additive semigroup of exponential weights\\
Binomial-version $q$-Fuss--Catalan & $(1-Q)\Pinf^{-1}q^{r n(n-1)}$ & Powers of $q^{-s}$; explicit cancellation rules\\
$\lvert\Sp(2n,q)\rvert$ & $\mathcal P_{Q^2}q^{2n^2+n}$ & Only even weights in $q^{-s}$\\
Unitary isometry groups & $\mathcal P_{-Q}q^{n^2}$ & Two parity sheets with different first corrections\\
Generalized Galois numbers, $r$ blocks & $\T^{(r)}_{n\bmod r}\Pinf^{1-r}q^{(r-1)n^2/(2r)}$ & Root-lattice theta coefficients; powers of $q^{-s/r}$\\
Normalized product limits & A nonzero finite limit & A logarithm plus a convergent power or Puiseux series in the deficit\\
Primitive necklaces, irreducible polynomials & $q^n/n$ & $W_{-1}$ center; arithmetic sectors or fixed-radical analytic sheets\\
\bottomrule
\end{longtable}

The finite product definitions and $q$-binomial identities used below are standard; DLMF and the Wolfram documentation specify their conventions \cite{dlmfq,wolframqbin,wolframpoch}. OEIS is used to identify representative integer specializations, not as a substitute for proof. The leading theta asymptotics of generalized Galois numbers have an established literature \cite{kousidis,bliem}; our development gives a direct coefficient-level construction and then carries it through inversion.

\section{What is being inverted?}
\label{sec:interpolation}

\subsection{A sequence and an interpolation are different objects}

Let $a_n$ be eventually increasing. Its integer threshold is an intrinsically defined function,
\begin{equation}\label{eq:threshold}
 N_a(Y)=\min\{n\ge n_0:a_n\ge Y\}.
\end{equation}
A continuous inverse instead requires a function $F(x)$ satisfying $F(n)=a_n$ on the relevant integers. Unless an interpolation has been specified, exponentially small terms of $F^{-1}$ are not determined by the sequence alone.

\begin{proposition}[Interpolation can change a flat inverse term]\label{prop:ambiguity}
Suppose $F>0$ is an eventually increasing smooth interpolation, and its logarithmic derivative is eventually bounded below by a positive constant. For fixed $\sigma>0$ and sufficiently small real $\eta$, set
\[
 \widetilde F(x)=F(x)\exp\!\bigl(\eta\ee^{-\sigma x}\sin\pi x\bigr).
\]
Then $\widetilde F(n)=F(n)$ for every integer $n$, and $\widetilde F$ is eventually increasing. If $x=F^{-1}(Y)$, then, locally as $\eta\to0$,
\begin{equation}\label{eq:interpolation-shift}
 \widetilde F^{-1}(Y)-x
 =-\eta\frac{\ee^{-\sigma x}\sin\pi x}{(\log F)'(x)}+O(\eta^2).
\end{equation}
\end{proposition}
\begin{proof}
The sine factor vanishes at integers. The derivative of the added logarithmic term is $O(\ee^{-\sigma x})$, so eventual positivity of the original logarithmic derivative is preserved. Differentiate
$\log F(x(\eta))+\eta\ee^{-\sigma x(\eta)}\sin\pi x(\eta)=\log Y$
with respect to $\eta$ at zero. The implicit function theorem gives the stated expansion.
\end{proof}

In the product examples we therefore \emph{define} the interpolation by a ratio of convergent infinite products. In the theta examples we define one analytic function for each residue class. In the arithmetic examples we distinguish exact integer formulas from explicitly chosen functions on fixed-radical subsequences.

\subsection{Convergent exponential expansions are transseries too}

Consider
\begin{equation}\label{eq:basic-shape}
 F(x)=\exp\bigl(P(x)+B(\ee^{-\lambda x})\bigr),\qquad
 B(t)=\sum_{m\ge1}b_m t^m,
\end{equation}
where $B$ is analytic at zero and $\lambda>0$. Its relative expansion is
\[
 F(x)=\ee^{P(x)}\sum_{m\ge0}v_m\ee^{-m\lambda x},\qquad v_0=1.
\]
For every fixed $m$, the factor $\ee^{-m\lambda x}$ is smaller than every negative power of $x$. Thus these are exponentially separated transseries sectors even when the series in the small variable $t$ converges. There is no requirement that a transseries be divergent.

The inverse will be written
\begin{equation}\label{eq:inverse-shape-intro}
 F^{-1}(Y)=s+\sum_{m\ge1}d_m(s)\ee^{-m\lambda s},\qquad
 P(s)=\log Y.
\end{equation}
For a quadratic $P$, the $d_m$ are finite polynomials in reciprocal powers of $P'(s)$, not merely formal infinite expansions inside each sector. Keeping the full center $s$ is important: a missing constant in $P$ creates an error of order $1/P'(s)$, which is much larger than any of the exponential corrections.

\subsection{Rounding belongs after inversion}

For an increasing interpolation valid on all sufficiently large integers,
\begin{equation}\label{eq:rounding}
 N_a(Y)=\lceil F^{-1}(Y)\rceil
\end{equation}
when the threshold lies in that range. For a sheet matching $n\equiv\eps\pmod r$, the corresponding threshold is
\begin{equation}\label{eq:residue-rounding}
 N_{\eps}(Y)=\eps+r\left\lceil\frac{F_{\eps}^{-1}(Y)-\eps}{r}\right\rceil.
\end{equation}
The overall threshold is the minimum of these sheet thresholds, after any finite initial range has been checked. A tiny real error does not automatically certify a ceiling when the true inverse is near an integer. Section~\ref{sec:certification} gives residual and rounding criteria.

\section{The analytic tail of a finite \texorpdfstring{$q$}{q}-product}
\label{sec:qtail}

\subsection{The common building block}

Define
\begin{equation}\label{eq:pinf}
 \Pinf=(Q;Q)_\infty=\prod_{j=1}^{\infty}(1-Q^j)>0,
 \qquad
 \A(t)=-\log(Qt;Q)_\infty.
\end{equation}
The logarithm is the analytic branch equal to zero at $t=0$. The first possible singularity is at $t=Q^{-1}$.

\begin{lemma}[Tail coefficients and an explicit remainder]\label{lem:Atail}
For $|t|<Q^{-1}$,
\begin{equation}\label{eq:Aseries}
 \A(t)=\sum_{m\ge1}\alpha_m t^m,
 \qquad \alpha_m=\frac{Q^m}{m(1-Q^m)}.
\end{equation}
For $M\ge0$ and $Q|t|<1$,
\begin{equation}\label{eq:Atailbound}
 \left|\A(t)-\sum_{m=1}^M\alpha_m t^m\right|
 \le \frac{(Q|t|)^{M+1}}
 {(M+1)(1-Q)(1-Q|t|)}.
\end{equation}
\end{lemma}
\begin{proof}
Expand $-\log(1-Q^jt)=\sum_{m\ge1}Q^{jm}t^m/m$ and sum over $j\ge1$. Absolute convergence on compact subsets of $|t|<Q^{-1}$ permits interchanging the sums and gives \eqref{eq:Aseries}. For the remainder use $1-Q^m\ge1-Q$, $m\ge M+1$, and then sum a geometric series.
\end{proof}

The interpolation of a finite product that will be used throughout is
\begin{equation}\label{eq:product-interpolation}
 \mathcal P_Q(x)=\frac{(Q;Q)_\infty}{(Q^{x+1};Q)_\infty}
 =\Pinf\exp\A(Q^x),\qquad x>-1.
\end{equation}
At $x=n\in\N$ it is precisely $\pp{n}$. It decreases to $\Pinf$ as $x\to\infty$, since $Q^x$ decreases and $\A$ has positive coefficients. The notation $\mathcal P_Q(x)$ is distinguished from the limiting constant $\Pinf$.

\subsection{Forward coefficients without logarithms}

Two Euler expansions, in the convention of \cite{dlmfq}, are
\begin{align}
 (Qt;Q)_\infty
 &=\sum_{j\ge0}\frac{(-1)^jQ^{j(j+1)/2}}{\pp{j}}t^j,
 \label{eq:euler-entire}\\
 \frac1{(Qt;Q)_\infty}
 &=\sum_{j\ge0}\frac{Q^j}{\pp{j}}t^j,
 \qquad |Qt|<1.
 \label{eq:euler-recip}
\end{align}
For completeness, write $E(z)=(z;Q)_\infty$. Its functional equation
$E(z)=(1-z)E(Qz)$ gives
$c_j=-Q^{j-1}c_{j-1}/(1-Q^j)$ for its coefficients. This proves the first expansion, including entire convergence. The reciprocal equation
$(1-z)/E(z)=1/E(Qz)$ gives $c_j=c_{j-1}/(1-Q^j)$ for $1/E(z)$ and proves the second. In particular,
\begin{equation}\label{eq:P-forward}
 \mathcal P_Q(x)=\Pinf\sum_{m\ge0}\frac{Q^{m(x+1)}}{\pp{m}}.
\end{equation}

More generally, if $B(t)=\sum b_m t^m$, then
\begin{equation}\label{eq:forward-Bell}
 \coef{t}{m}\ee^{B(t)}
 =\sum_{\substack{k_1,\ldots,k_m\ge0\\\sum j k_j=m}}
       \prod_{j=1}^m\frac{b_j^{k_j}}{k_j!}.
\end{equation}
This is a finite formula at every weight. It follows by multiplying
$\prod_j\sum_{k\ge0}b_j^k t^{jk}/k!$; only $j\le m$ can contribute to weight $m$.

\subsection{A finite-product coefficient rule}

Many families can be reduced immediately to this format. Suppose
\begin{equation}\label{eq:product-rule}
 \ee^{B(t)}=\prod_{\nu=1}^J
       (Q^{\beta_\nu}t^{r_\nu};Q^{e_\nu})_\infty^{\eta_\nu},
 \qquad r_\nu\in\N_{>0},\quad e_\nu>0,
\end{equation}
where the powers use the branch equal to one at $t=0$. Real $\beta_\nu$ and real or complex $\eta_\nu$ are permitted. Taking the logarithm gives the explicit formula
\begin{equation}\label{eq:product-log-coeff}
 b_m=-\sum_{\substack{1\le\nu\le J\\r_\nu\mid m}}
 \eta_\nu\frac{Q^{\beta_\nu m/r_\nu}}
 {(m/r_\nu)(1-Q^{e_\nu m/r_\nu})}.
\end{equation}
The identity is valid in a neighborhood of zero where all the factors are nonzero. Combining \eqref{eq:product-log-coeff}, \eqref{eq:forward-Bell}, and the inversion theorem below gives an explicit procedure for the forward and inverse coefficients of any such finite product.

\section{A convergent inversion theorem with finite coefficient formulas}
\label{sec:engine}

\subsection{Statement and analytic control}

Let
\begin{equation}\label{eq:Pquadratic}
 P(x)=a x^2+b x+c,
\end{equation}
where either $a>0$, or $a=0$ and $b>0$. The function $B$ in \eqref{eq:basic-shape} is fixed, analytic at zero, and real for sufficiently small real arguments. Let $L=\log Y$ and choose the large real solution of $P(s)=L$:
\begin{equation}\label{eq:center}
 s=\begin{cases}
 \displaystyle\frac{-b+\sqrt{b^2+4a(L-c)}}{2a},&a>0,\\[0.8em]
 \displaystyle\frac{L-c}{b},&a=0.
 \end{cases}
 \qquad D=P'(s)=2as+b.
\end{equation}

\begin{theorem}[Convergent inverse transseries]\label{thm:inverse}
For all sufficiently large real $Y$, the function
$F(x)=\exp(P(x)+B(\ee^{-\lambda x}))$ is increasing near its large real inverse, and
\begin{equation}\label{eq:inverse-full}
 F^{-1}(Y)=s+\sum_{m\ge1}d_m(D)\,t_0^m,
 \qquad t_0=\ee^{-\lambda s}.
\end{equation}
The series converges. In fact, there exist positive constants $R$, $C$, and $s_0$, independent of $s\ge s_0$, such that
\begin{equation}\label{eq:inverse-remainder}
 \left|F^{-1}(Y)-s-\sum_{m=1}^M d_m(D)t_0^m\right|
 \le \frac{C}{D}\frac{(t_0/R)^{M+1}}{1-t_0/R}
 \quad(0<t_0<R).
\end{equation}
Each $d_m$ is an explicitly computable finite polynomial in $D^{-1}$, with coefficients polynomial in $a,\lambda,b_1,\ldots,b_m$. The highest possible power of $D^{-1}$ is $2m-1$.
\end{theorem}

\begin{proof}
Write $x=s+\delta$. The inverse equation is exactly
\begin{equation}\label{eq:implicit-delta}
 D\delta+a\delta^2+B(t\ee^{-\lambda\delta})=0,
\end{equation}
with $t=t_0$ at the end. Choose $D_0>0$ such that $D\ge D_0$ in the range considered. Choose $\rho>0$ so small that $a\rho\le D_0/4$; this restriction is unnecessary when $a=0$. Then choose $R>0$ so small that $B$ is analytic on the closed disk $|u|\le R\ee^{\lambda\rho}$ and
\[
 \sup_{|u|\le R\ee^{\lambda\rho}}|B(u)|<D_0\rho/4.
\]
On $|\delta|=\rho$ and $|t|\le R$, the sum of the last two terms of \eqref{eq:implicit-delta} has modulus less than $D|\delta|$. Rouch\'e's theorem gives exactly one zero in the disk, counted with multiplicity. The zero is therefore simple and depends holomorphically on $t$. It is the branch satisfying $\delta(0)=0$.

Since $B(0)=0$, there is a constant $K$ such that
$|B(t\ee^{-\lambda\delta})|\le K|t|$ on these disks. Equation \eqref{eq:implicit-delta} implies
\[
 (D-a\rho)|\delta|\le K|t|,
 \qquad |\delta|\le \frac{4K|t|}{3D}.
\]
Cauchy's coefficient estimates on a slightly smaller fixed $t$-disk yield
$|d_m|\le C D^{-1}R^{-m}$ after changing $C$ and $R$. Summing the resulting geometric tail proves \eqref{eq:inverse-remainder}. For large $s$, $t_0<R$.

Finally,
\[
 (\log F)'(x)=2ax+b-\lambda\ee^{-\lambda x}B'(\ee^{-\lambda x})>0
\]
for large $x$, so the local branch is the large real inverse. The finite polynomial assertion follows from the explicit formula proved next.
\end{proof}

\subsection{A Lagrange--B\"urmann formula}

Define the local response to a prescribed logarithmic perturbation by
\begin{equation}\label{eq:phi}
 \Phi_D(u)=\begin{cases}
 \displaystyle\frac{-D+\sqrt{D^2-4aB(u)}}{2a},&a>0,\\[0.8em]
 -B(u)/D,&a=0,
 \end{cases}
\end{equation}
where the square root equals $D$ at zero. Then $D\Phi_D+a\Phi_D^2=-B(u)$.

\begin{theorem}[Closed coefficient extraction]\label{thm:coefficient}
For $m\ge1$,
\begin{equation}\label{eq:inverse-extraction}
 \boxed{\displaystyle
 d_m(D)=-\frac1{m\lambda}\coef{u}{m}
          \exp\bigl(-m\lambda\Phi_D(u)\bigr).}
\end{equation}
Equivalently, put
\begin{equation}\label{eq:Cmk}
 C_{m,k}=\coef{u}{m}B(u)^k
 =\sum_{\substack{j_1,\ldots,j_k\ge1\\j_1+\cdots+j_k=m}}
       b_{j_1}\cdots b_{j_k}.
\end{equation}
Then the completely finite, nonrecursive formula is
\begin{equation}\label{eq:finite-dm}
 \boxed{\displaystyle
 d_m=-\sum_{k=1}^m\frac{C_{m,k}}{k}
 \sum_{j=0}^{k-1}\binom{k+j-1}{j}
 \frac{a^j(m\lambda)^{k-1-j}}
 {(k-1-j)!\,D^{k+j}}.}
\end{equation}
The $a=0$ case is obtained by retaining only $j=0$.
\end{theorem}

\begin{proof}
Set $u=t\ee^{-\lambda\delta}$. Equation \eqref{eq:implicit-delta} becomes
$\delta=\Phi_D(u)$ and $u=t\exp(-\lambda\Phi_D(u))$.
The elementary Lagrange--B\"urmann identity gives
\[
 [t^m]\Phi_D(u(t))
 =\frac1m[u^{m-1}]\Phi_D'(u)\ee^{-m\lambda\Phi_D(u)}.
\]
Differentiating the exponential and using
$[u^{m-1}]H'(u)=m[u^m]H(u)$ proves \eqref{eq:inverse-extraction}.

To obtain a finite formula, write $v=B(u)$ and
$w(v)=-\Phi_D(u)$. Thus $Dw-aw^2=v$, or
$w=v/(D-aw)$. A second application of Lagrange inversion gives, for $k\ge1$,
\[
 [v^k]\ee^{m\lambda w(v)}
 =\frac{m\lambda}{k}[w^{k-1}]
      \ee^{m\lambda w}(D-aw)^{-k}.
\]
Expand the exponential and the negative binomial power. Their coefficient of $w^{k-1}$ is
\[
 \sum_{j=0}^{k-1}\binom{k+j-1}{j}
 \frac{a^j(m\lambda)^{k-1-j}}{(k-1-j)!D^{k+j}}.
\]
Substitute this expression into \eqref{eq:inverse-extraction}, and expand $B(u)^k$. Only $k\le m$ contributes. This proves \eqref{eq:finite-dm} and also the bound $k+j\le2m-1$ on the reciprocal power of $D$.
\end{proof}

\subsection{The first three sectors and a computation recurrence}

The first coefficients are
\begin{align}
 d_1={}&-\frac{b_1}{D},\label{eq:d1}\\
 d_2={}&-\frac{b_2}{D}-\frac{\lambda b_1^2}{D^2}
                   -\frac{a b_1^2}{D^3},\label{eq:d2}\\
 d_3={}&-\frac{b_3}{D}-\frac{3\lambda b_1b_2}{D^2}
 -\frac{2a b_1b_2+\tfrac32\lambda^2b_1^3}{D^3}
 -\frac{3a\lambda b_1^3}{D^4}
 -\frac{2a^2b_1^3}{D^5}.\label{eq:d3}
\end{align}
The $D^{-3}$ term in $d_2$ comes from curvature of the dominant phase, while its $D^{-2}$ term comes from shifting the exponential argument. Keeping only the response $-B(t_0)/D$ misses both.

For computation, let $\delta_{<m}(t)=\sum_{j=1}^{m-1}d_jt^j$. Coefficient comparison in \eqref{eq:implicit-delta} gives the triangular recurrence
\begin{equation}\label{eq:dm-recurrence}
 d_m=-\frac1D\left(
 a\sum_{i=1}^{m-1}d_i d_{m-i}
 +\sum_{k=1}^m b_k\coef{t}{m-k}
     \ee^{-k\lambda\delta_{<m}(t)}\right).
\end{equation}
This is convenient for numerical work; \eqref{eq:finite-dm} shows that no recurrence is required in the mathematical specification of a coefficient.

\begin{corollary}[Allowed weights]\label{cor:support}
Let $S=\{m\ge1:b_m\ne0\}$. A coefficient $d_m$ can be nonzero only when $m$ is a finite sum of elements of $S$. In particular, if $B(t)=\widehat B(t^r)$, then $d_m=0$ unless $r\mid m$.
\end{corollary}
\begin{proof}
In \eqref{eq:finite-dm}, every term contributing to $C_{m,k}$ has $m=j_1+\cdots+j_k$ with $j_i\in S$. This proves the assertion. Some allowed weights may still cancel; membership in the semigroup is not a guarantee of nonvanishing.
\end{proof}

\begin{corollary}[Direction of the first correction]\label{cor:sign}
If $b_j=0$ for $j<r$ and $b_r\ne0$, then
\[
 F^{-1}(Y)-s=-\frac{b_r}{D}\ee^{-r\lambda s}
 +O\!\left(\frac{\ee^{-(r+1)\lambda s}}{D}\right).
\]
If all $b_m\ge0$, then all $d_m\le0$ for $a\ge0$.
\end{corollary}
\begin{proof}
The first assertion follows from the finite formula and the remainder bound. In the second assertion all factors in \eqref{eq:finite-dm}, except the initial minus sign, are nonnegative.
\end{proof}

\section{General linear groups, complete flags, and \texorpdfstring{$q$}{q}-factorials}
\label{sec:GL}

\subsection{Ordered bases and general linear groups}

An invertible $n\times n$ matrix over $\F_q$ is an ordered basis. Its first column can be any nonzero vector, its second can be any vector outside the first column's span, and so on. Therefore
\begin{equation}\label{eq:GLcount}
 g_n(q)=|\GL(n,q)|=\prod_{j=0}^{n-1}(q^n-q^j)
       =q^{n^2}\pp{n}.
\end{equation}
The specialization $q=2$ is OEIS A002884 \cite{oeisGL}. Our chosen interpolation is
\begin{equation}\label{eq:GLinterp}
 g_q(x)=\Pinf\exp\bigl(hx^2+\A(Q^x)\bigr).
\end{equation}
It agrees with \eqref{eq:GLcount} at every nonnegative integer. Its exact forward transseries is especially simple:
\begin{equation}\label{eq:GLforward}
 g_q(x)=\Pinf q^{x^2}
       \sum_{m\ge0}\frac{Q^{m(x+1)}}{\pp{m}}.
\end{equation}
For $q\ge2$, the interpolation is already increasing for $x\ge1$. Indeed,
\[
 \frac{(\log g_q)'(x)}h
 =2x-\sum_{m\ge1}\frac{Q^{m(x+1)}}{1-Q^m},
\]
and the sum is bounded by
$Q^{x+1}/((1-Q)(1-Q^{x+1}))\le2/3$ in this range.
For arbitrary fixed $q>1$, eventual monotonicity is sufficient and follows from Theorem~\ref{thm:inverse}.

Put
\begin{equation}\label{eq:GLcenter}
 s=\sqrt{\frac{\log(Y/\Pinf)}h},\qquad D=2hs,\qquad t_0=Q^s.
\end{equation}
Using $a=\lambda=h$ and $b_m=\alpha_m$ gives
\begin{align}\label{eq:GLinverse}
 g_q^{-1}(Y)=s
 &-\frac{\alpha_1}{D}t_0\\[-0.4em]
 &-\left(\frac{\alpha_2}{D}
       +\frac{h\alpha_1^2}{D^2}
       +\frac{h\alpha_1^2}{D^3}\right)t_0^2
 +O(t_0^3/D).\notag
\end{align}
Every later coefficient is given by \eqref{eq:finite-dm}. All corrections are nonpositive. This agrees with the elementary inequality
$g_q(x)>\Pinf q^{x^2}$, which puts the true inverse below the leading center.

For $q=2$,
\[
 \alpha_1=1,\qquad\alpha_2=\frac16,\qquad
 \alpha_3=\frac1{21},\qquad\alpha_4=\frac1{60}.
\]
Thus even the first two inverse sectors can be evaluated with no special functions beyond a logarithm and a square root, once $\Pinf$ is known.

\subsection{Complete flags and the inversion statistic on permutations}

A complete flag in $\F_q^n$ is a chain
$0=V_0\subset V_1\subset\cdots\subset V_n=\F_q^n$ with $\dim V_j=j$.
Each flag has exactly
$(q-1)^n q^{n(n-1)/2}$ ordered bases adapted to it: the $j$th vector may be any element of $V_j\setminus V_{j-1}$.
Dividing \eqref{eq:GLcount} by this number gives
\begin{equation}\label{eq:qfactorial}
 [n]_q!=\prod_{j=1}^n\frac{q^j-1}{q-1}
 =q^{n(n+1)/2}(q-1)^{-n}\pp{n}.
\end{equation}
At $q=2$ these are the products $\prod_{j=1}^n(2^j-1)$ in A005329 \cite{oeisflags}.

There is also a weighted permutation interpretation. If $\operatorname{inv}(\pi)$ denotes the inversion number of a permutation, insertion of the largest element creates any one of $0,1,\ldots,n-1$ new inversions. Induction therefore proves
\begin{equation}\label{eq:Mahonian}
 \sum_{\pi\in S_n}q^{\operatorname{inv}(\pi)}=[n]_q!.
\end{equation}
This relates a finite-field count to an ordinary combinatorial generating polynomial.

For the analytic interpolation, use
\begin{equation}\label{eq:flagphase}
 a=\frac h2,\quad b=\frac h2-\log(q-1),\quad
 c=\log\Pinf,\quad B(t)=\A(t),\quad\lambda=h.
\end{equation}
The center is the quadratic root \eqref{eq:center}, and all inverse coefficients follow from the same $\alpha_m$ used for general linear groups. Only $a$ and $D$ change. In particular,
\begin{equation}\label{eq:flaginverse}
 ([x]_q!)^{-1}(Y)
 =s-\frac{\alpha_1}{D}Q^s
 -\left(\frac{\alpha_2}{D}
   +\frac{h\alpha_1^2}{D^2}
   +\frac{h\alpha_1^2}{2D^3}\right)Q^{2s}
 +O(Q^{3s}/D).
\end{equation}
Here the notation on the left means the inverse of the interpolation just defined, not the reciprocal of a factorial.

\subsection{The other side of \texorpdfstring{$q=1$}{q = 1}: an exponential rather than quadratic phase}
\label{sec:qfactorial-smallbase}

Use $Q\in(0,1)$ as the base of the usual $q$-factorial. Then
\begin{equation}\label{eq:smallbase-factorial}
 [n]_Q!=\frac{\pp{n}}{(1-Q)^n},\qquad
 f_Q(x)=\Pinf\exp\bigl(\beta x+\A(Q^x)\bigr),
 \quad \beta=-\log(1-Q)>0.
\end{equation}
The logarithmic phase is now linear. Set
\begin{equation}\label{eq:smallbase-center}
 s=\frac{\log(Y/\Pinf)}\beta,\qquad
 \rho=\frac h\beta,\qquad
 t_0=Q^s=(Y/\Pinf)^{-\rho}.
\end{equation}
The inverse is a convergent expansion in a generally nonintegral power of $Y$:
\begin{align}\label{eq:smallbase-inverse}
 f_Q^{-1}(Y)=s
 -\frac{\alpha_1}{\beta}t_0
 -\left(\frac{\alpha_2}{\beta}
       +\frac{h\alpha_1^2}{\beta^2}\right)t_0^2
 +O(t_0^3).
\end{align}
The exact all-order formula simplifies to
\begin{equation}\label{eq:linear-extraction}
 d_m=-\frac1{mh}\coef{u}{m}\exp\left(\frac{mh}{\beta}\A(u)\right).
\end{equation}

Consequently the same polynomial sequence $[n]_q!$, viewed at a fixed real base, has three different large-index regimes. For $q>1$ its logarithm is quadratic in $n$; for $0<q<1$ it is linear; and at $q=1$ it is $\log(n!)$, with dominant term $n\log n$. This is a genuine change of asymptotic scaling, not something obtained by substituting $q=1$ into a fixed-$q$ inverse transseries. Section~\ref{sec:double} studies a related simultaneous limit in detail.

\section{Gaussian binomial coefficients and prescribed partial flags}
\label{sec:gaussian}

\subsection{Counting subspaces}

For integers $0\le k\le n$, the number of $k$-dimensional subspaces of $\F_q^n$ is
\begin{equation}\label{eq:qbin-count}
 \qbin{n}{k}
 =\frac{\prod_{j=0}^{k-1}(q^n-q^j)}
        {\prod_{j=0}^{k-1}(q^k-q^j)}
 =q^{k(n-k)}\frac{\pp{n}}{\pp{k}\pp{n-k}}.
\end{equation}
The first quotient counts linearly independent ordered $k$-tuples and divides by the number of bases of each subspace. This proves both the enumerative interpretation and the product identity. The definition agrees with \texttt{QBinomial[n,k,q]} \cite{wolframqbin}.

\subsection{Fixed rank: algebraic inverse series}
\label{sec:fixedrank}

Fix $k\ge1$. The finite-product interpolation is
\begin{equation}\label{eq:fixedrank-interp}
 H_{k,q}(x)=\frac{q^{k(x-k)}}{\pp{k}}
              \prod_{j=0}^{k-1}(1-q^jQ^x),\qquad x>k-1.
\end{equation}
It is positive and increasing in this range, as is evident from the first quotient in \eqref{eq:qbin-count} with $n$ replaced by $x$. In the notation of the inversion theorem,
\begin{align}
 P(x)&=khx-hk^2-\log\pp{k},\label{eq:fixedrank-phase}\\
 B(t)&=\sum_{j=0}^{k-1}\log(1-q^jt),\qquad
 b_m=-\frac{q^{km}-1}{m(q^m-1)}.\label{eq:fixedrank-tail}
\end{align}
The center and small parameter are
\begin{equation}\label{eq:fixedrank-center}
 s=\frac{\log Y+hk^2+\log\pp{k}}{kh},
 \qquad t_0=q^{-k}\bigl(\pp{k}Y\bigr)^{-1/k}.
\end{equation}
Thus a convergent inverse Puiseux expansion follows:
\begin{equation}\label{eq:fixedrank-inverse}
 H_{k,q}^{-1}(Y)
 =s+\sum_{m\ge1}d_m t_0^m,
 \qquad
 d_m=-\frac1{mh}\coef{u}{m}
       \prod_{j=0}^{k-1}(1-q^ju)^{m/k}.
\end{equation}
All powers are expanded at $u=0$ with constant term one. In particular,
\begin{equation}\label{eq:fixedrank-first}
 H_{k,q}^{-1}(Y)=s+\frac{[k]_q}{kh}\,t_0+O(t_0^2).
\end{equation}
Although the leading inverse contains a logarithm, its correction coefficients are constants: $D=kh$ is fixed.

There is an independent algebraic explanation of convergence. Writing $z=q^x$, the equation is
\begin{equation}\label{eq:fixedrank-algebraic}
 \prod_{j=0}^{k-1}(z-q^j)=Y\prod_{j=0}^{k-1}(q^k-q^j).
\end{equation}
After scaling $z$ by $Y^{1/k}$, the desired large positive root is analytic in $Y^{-1/k}$ near zero, because the limiting polynomial has a simple positive root. Taking a logarithm gives \eqref{eq:fixedrank-inverse}.

For $k=1$ this reduces to
\[
 H_{1,q}^{-1}(Y)=\frac{\log(1+(q-1)Y)}h.
\]
For $k=2$, put $K_q=(q^2-1)(q^2-q)$. The inverse is exactly
\begin{equation}\label{eq:k2-inverse}
 H_{2,q}^{-1}(Y)=\frac1h\log\left(
 \frac{q+1+\sqrt{(q-1)^2+4K_qY}}2\right).
\end{equation}
These elementary cases provide useful checks on the coefficient formula.

\subsection{Proportionally scaled rank: a quadratic inverse center}

Let $u,v$ be fixed positive integers. Define
\begin{equation}\label{eq:scaled-gaussian}
 H_{u,v}(x)=\frac{q^{uvx^2}}{\Pinf}
 \exp\bigl(\A(t^{u+v})-\A(t^u)-\A(t^v)\bigr),
 \qquad t=Q^x.
\end{equation}
At integers $x=n$ it is $\qbin{(u+v)n}{un}$. The logarithmic coefficients are
\begin{equation}\label{eq:scaled-gaussian-b}
 b_m=\one_{u+v\mid m}\alpha_{m/(u+v)}
     -\one_{u\mid m}\alpha_{m/u}
     -\one_{v\mid m}\alpha_{m/v}.
\end{equation}
Use $a=uvh$, $b=0$, $c=-\log\Pinf$, and $\lambda=h$ in the inversion theorem. The first exponential weight is $\min(u,v)$; if $u=v$, its coefficient is doubled. The available inverse weights belong to the additive semigroup generated by the nonzero degrees in \eqref{eq:scaled-gaussian-b}.

For the central case $u=v=1$,
\begin{equation}\label{eq:central-B}
 B(t)=\A(t^2)-2\A(t),\qquad
 b_m=-2\alpha_m+\one_{2\mid m}\alpha_{m/2}.
\end{equation}
The corresponding $q=2$ sequence is A006098 \cite{oeiscentral}. Its center is
\begin{equation}\label{eq:central-center}
 s=\sqrt{\frac{\log(Y\Pinf)}h},\qquad D=2hs.
\end{equation}
The first two inverse sectors are
\begin{align}\label{eq:central-inverse}
 H_{1,1}^{-1}(Y)=s
 &+\frac{2\alpha_1}{D}Q^s\\[-0.3em]
 &-\left(\frac{\alpha_1-2\alpha_2}{D}
       +\frac{4h\alpha_1^2}{D^2}
       +\frac{4h\alpha_1^2}{D^3}\right)Q^{2s}
 +O(Q^{3s}/D).\notag
\end{align}
The leading correction has the opposite sign from that for $\GL(n,q)$.

An explicit forward coefficient formula follows directly from the Euler expansions:
\begin{equation}\label{eq:central-forward}
 \ee^{B(t)}=\frac{(Qt;Q)_\infty^2}{(Qt^2;Q)_\infty},\qquad
 [t^m]\ee^{B(t)}=
 \sum_{r+s+2j=m}
 \frac{(-1)^{r+s}Q^{r(r+1)/2+s(s+1)/2+j}}
 {\pp{r}\pp{s}\pp{j}}.
\end{equation}
The sum ranges over nonnegative integers and is finite. Equations \eqref{eq:central-B} and \eqref{eq:finite-dm} give a similarly finite formula for every inverse coefficient.

\subsection{Prescribed partial flags}

Fix positive integers $r_1,\ldots,r_\ell$ and write
\[
 R=r_1+\cdots+r_\ell,\qquad
 A_0=\sum_{1\le i<j\le\ell}r_i r_j.
\]
A flag with successive quotient dimensions $r_1n,\ldots,r_\ell n$ is counted by the Gaussian multinomial coefficient
\begin{equation}\label{eq:multinomial}
 \genfrac{[}{]}{0pt}{}{Rn}{r_1n,\ldots,r_\ell n}_q
 =q^{A_0n^2}\frac{\pp{Rn}}{\prod_i\pp{r_i n}}.
\end{equation}
It follows either by successive subspace choices or by dividing adapted bases by the relevant block-upper-triangular stabilizer. The interpolation has
\begin{equation}\label{eq:partialflag-phase}
 P(x)=A_0hx^2+(1-\ell)\log\Pinf,
 \qquad
 B(t)=\A(t^R)-\sum_{i=1}^{\ell}\A(t^{r_i}).
\end{equation}
Thus
\begin{equation}\label{eq:partialflag-b}
 b_m=\one_{R\mid m}\alpha_{m/R}
     -\sum_{i:r_i\mid m}\alpha_{m/r_i}.
\end{equation}
All proportional partial-flag families are consequently covered by one inversion theorem and one explicit arithmetic coefficient sequence. Repeated block sizes automatically produce their correct multiplicities; they should not be reduced to distinct sizes before applying this formula.

\section{Binomial-version \texorpdfstring{$q$}{q}-Catalan and \texorpdfstring{$q$}{q}-Fuss--Catalan numbers}
\label{sec:catalan}

\subsection{Definition and polynomiality}

There are several inequivalent $q$-Catalan conventions. Here we use only the binomial quotient
\begin{equation}\label{eq:qfuss-definition}
 C_{r,n}(q)=\frac1{[rn+1]_q}\qbin{(r+1)n}{n},
 \qquad r\in\N_{>0},\qquad
 [m]_q=\frac{q^m-1}{q-1}.
\end{equation}
At $r=1$ and $q=2$ this is A015030 \cite{oeiscatalan}. At $q=1$, by continuity, it is the usual Fuss--Catalan number
$\binom{(r+1)n}{n}/(rn+1)$. No identification with a different area-weighted or recursively defined $q$-Catalan polynomial is intended.

\begin{proposition}\label{prop:qfuss-polynomial}
For positive integers $r,n$, the expression \eqref{eq:qfuss-definition} is a polynomial in $q$ with integer coefficients.
\end{proposition}
\begin{proof}
The exponent of the cyclotomic polynomial $\Phi_d(q)$ in
$\qbin{(r+1)n}{n}$ is
\[
 e_d=\left\lfloor\frac{(r+1)n}{d}\right\rfloor
     -\left\lfloor\frac n d\right\rfloor
     -\left\lfloor\frac{rn}{d}\right\rfloor\in\{0,1\}.
\]
This follows by factoring each $q^j-1$ into cyclotomic factors. A factor of $[rn+1]_q$ has $d>1$ and $d\mid rn+1$. Write $rn=ad-1$. Then
$e_d=1+\lfloor(n-1)/d\rfloor-\lfloor n/d\rfloor=1$, because $d\nmid n$ follows from $\gcd(n,rn+1)=1$. Hence every denominator factor cancels, and all remaining cyclotomic exponents are nonnegative integers.
\end{proof}
This proof establishes integer coefficients. No coefficient-positivity assertion is needed in the analysis below; positivity of the quotient for real $q>1$ follows directly from its factors.

\subsection{An exact logarithmic tail}

From \eqref{eq:scaled-gaussian}, and
$[rx+1]_q=q^{rx+1}(1-Qt^r)/(q-1)$, obtain the interpolation
\begin{equation}\label{eq:qfuss-interp}
 C_{r,q}(x)=\frac{1-Q}{\Pinf}\,q^{r x(x-1)}\ee^{B_r(t)},
 \qquad t=Q^x,
\end{equation}
where
\begin{equation}\label{eq:qfuss-B}
 B_r(t)=\A(t^{r+1})-\A(t)-\A(t^r)-\log(1-Qt^r).
\end{equation}
Consequently
\begin{equation}\label{eq:qfuss-b}
 b_m=\one_{r+1\mid m}\alpha_{m/(r+1)}-\alpha_m
 -\one_{r\mid m}\alpha_{m/r}
 +\one_{r\mid m}\frac{Q^{m/r}}{m/r}.
\end{equation}
This formula includes $r=1$ correctly; two copies of $-\A(t)$ then occur.

For $r=1$, the leading asymptotic specializes to
\begin{equation}\label{eq:qCatalan-leading}
 C_{1,n}(2)\sim\frac1{\mathcal P_{1/2}}\,2^{n^2-n-1}.
\end{equation}
This is the asymptotic constant recorded by Vladimir Reshetnikov in the 2021 contribution to A015030 \cite{oeiscatalan}. Formula \eqref{eq:qfuss-b} supplies all of its exponential corrections, while Theorem~\ref{thm:inverse} supplies all inverse corrections.

\subsection{Inverse formulas}

The exact quadratic center is
\begin{equation}\label{eq:qfuss-center}
 s=\frac12+\sqrt{\frac14+
       \frac{\log\bigl(Y\Pinf/(1-Q)\bigr)}{rh}},
 \qquad D=rh(2s-1).
\end{equation}
Use $a=rh$ and $\lambda=h$ in \eqref{eq:finite-dm}. For $r=1$,
\begin{equation}\label{eq:qCatalan-b12}
 b_1=-\frac{Q(1+Q)}{1-Q},\qquad
 b_2=\alpha_1-2\alpha_2+\frac{Q^2}{2}.
\end{equation}
Hence, with $\gamma=Q(1+Q)/(1-Q)$,
\begin{align}\label{eq:qCatalan-inverse}
 C_{1,q}^{-1}(Y)=s
 &+\frac{\gamma}{D}Q^s\\[-0.3em]
 &-\left(\frac{b_2}{D}
       +\frac{h\gamma^2}{D^2}
       +\frac{h\gamma^2}{D^3}\right)Q^{2s}
 +O(Q^{3s}/D).\notag
\end{align}
For $r>1$, $b_1=-\alpha_1$, so the first correction is
\begin{equation}\label{eq:qFuss-first}
 C_{r,q}^{-1}(Y)=s+\frac{\alpha_1}{rh(2s-1)}Q^s
                  +O(Q^{2s}/s).
\end{equation}
The case $r=1$ is therefore not obtained by treating all terms in \eqref{eq:qfuss-B} as having distinct weights: a collision occurs at the very first weight.

For example, at $q=2$, $r=1$, the coefficients are
\[
 b_1=-\frac32,\qquad b_2=\frac{19}{24},\qquad
 b_3=-\frac{3}{56}.
\]
The source and inverse coefficients are rational functions of $Q$, together with the powers of $h$ and $D$ dictated by the universal formula. There is no unresolved recurrence in these expressions.

\section{Symplectic and unitary orders: missing sectors and parity}
\label{sec:classical}

\subsection{Symplectic groups}

Let $V$ be a $2n$-dimensional symplectic vector space over $\F_q$. Choose the first vector $e_1\ne0$ of a symplectic basis in $q^{2n}-1$ ways. The equation $\langle e_1,f_1\rangle=1$ then has $q^{2n-1}$ solutions. The symplectic complement of their span has dimension $2n-2$. Iterating this argument proves
\begin{equation}\label{eq:Sp-count}
 s_n(q)=|\Sp(2n,q)|
 =q^{n^2}\prod_{j=1}^n(q^{2j}-1)
 =q^{2n^2+n}(Q^2;Q^2)_n.
\end{equation}
The $q=2$ specialization is A003923 \cite{oeisSp}.

Choose the interpolation
\begin{equation}\label{eq:Sp-interp}
 s_q(x)=\mathcal P_{Q^2}
   \exp\bigl(h(2x^2+x)+\mathcal A_{Q^2}(Q^{2x})\bigr).
\end{equation}
With $t=Q^x$, its logarithmic tail is
\begin{equation}\label{eq:Sp-B}
 B(t)=\mathcal A_{Q^2}(t^2),\qquad
 b_{2m}=\frac{Q^{2m}}{m(1-Q^{2m})},\qquad b_{2m-1}=0.
\end{equation}
The inverse therefore contains no odd weights in $Q^s$. Put
\begin{equation}\label{eq:Sp-center}
 s=\frac{-1+\sqrt{1+8\log(Y/\mathcal P_{Q^2})/h}}4,
 \qquad D=h(4s+1),\qquad \beta_m=\frac{Q^{2m}}{m(1-Q^{2m})}.
\end{equation}
Using the primitive small parameter $Q^{2s}$, so that $\lambda=2h$, gives
\begin{align}\label{eq:Sp-inverse}
 s_q^{-1}(Y)=s
 &-\frac{\beta_1}{D}Q^{2s}\\[-0.3em]
 &-\left(\frac{\beta_2}{D}
       +\frac{2h\beta_1^2}{D^2}
       +\frac{2h\beta_1^2}{D^3}\right)Q^{4s}
 +O(Q^{6s}/D).\notag
\end{align}
This illustrates why the exponential scale should be determined from the source coefficients, not guessed from a generic family resemblance.

\subsection{Unitary isometry groups and the order formula}

Let $U_n(q)$ denote the isometry group of a nondegenerate Hermitian form on $\F_{q^2}^n$. This is the full unitary isometry group, not the determinant-one subgroup and not a group of similitudes. The corresponding OEIS convention uses $\mathrm{GU}(n,q)$ \cite{oeisU}.

We include a count to fix all powers of $q$ and all constants. For the standard Hermitian form, write $N(v)=\sum_i v_i^{q+1}\in\F_q$. The norm map $\F_{q^2}^{\times}\to\F_q^{\times}$ has $q+1$ preimages for each nonzero value. If $\psi$ is a nontrivial additive character of $\F_q$, then
\[
 \sum_{z\in\F_{q^2}}\psi(z^{q+1})
 =1+(q+1)\sum_{a\in\F_q^{\times}}\psi(a)=-q.
\]
By character orthogonality, the number of vectors of any prescribed nonzero norm is
\begin{equation}\label{eq:unitnorm}
 \frac{q^{2n}-(-q)^n}{q}=q^{n-1}(q^n-(-1)^n).
\end{equation}
The orthogonal complement of a unit vector is again a nondegenerate Hermitian space of one smaller dimension. Counting orthonormal bases therefore yields
\begin{equation}\label{eq:U-count}
 u_n(q)=|U_n(q)|
 =q^{n(n-1)/2}\prod_{j=1}^n(q^j-(-1)^j)
 =q^{n^2}\prod_{j=1}^n(1-(-Q)^j).
\end{equation}
The reduction to the standard Hermitian form follows by the usual orthogonal basis construction; every nonzero diagonal norm can be rescaled to one because the norm map is onto.

\subsection{Two real analytic sheets}

Define the positive constant
\begin{equation}\label{eq:Pminus}
 \mathcal P_{-Q}=\prod_{j=1}^{\infty}(1-(-Q)^j)
 =(Q^2;Q^2)_\infty(-Q;Q^2)_\infty.
\end{equation}
For $\eps\in\{0,1\}$ define
\begin{align}
 B_{\eps}(t)&=\sum_{m\ge1}
 \frac{(-1)^{m(\eps+1)}Q^m}{m(1-(-Q)^m)}t^m,
 \label{eq:U-B}\\
 u_{q,\eps}(x)&=\mathcal P_{-Q}
                  \exp\bigl(hx^2+B_{\eps}(Q^x)\bigr).
 \label{eq:U-sheet}
\end{align}
The series in \eqref{eq:U-B} is analytic for $|t|<Q^{-1}$. At an integer $n\equiv\eps\pmod2$, the omitted tail of the finite product in \eqref{eq:U-count} starts with $(-Q)^{n+1}$; expansion of its logarithm gives exactly \eqref{eq:U-B}. Thus each sheet matches its specified parity.

Both inverse centers are
\begin{equation}\label{eq:U-center}
 s=\sqrt{\frac{\log(Y/\mathcal P_{-Q})}h},\qquad D=2hs.
\end{equation}
Since $b_1=(-1)^{\eps+1}Q/(1+Q)$,
\begin{equation}\label{eq:U-first}
 u_{q,\eps}^{-1}(Y)
 =s+(-1)^{\eps}\frac{Q}{(1+Q)D}Q^s+O(Q^{2s}/D).
\end{equation}
The even sheet lies above its center and the odd sheet below. Averaging the two sheets would erase the leading correction, although that correction is present on each integer subsequence.

At $q=2$,
\begin{equation}\label{eq:U-constant-check}
 \mathcal P_{-1/2}=1.210724130301059180136172856\ldots.
\end{equation}
Thus $u_n(2)\sim\mathcal P_{-1/2}2^{n^2}$. If this is written as
$c\,2^{n^2-1}$, the equivalent constant is
$c=2\mathcal P_{-1/2}=2.42144826060211836027\ldots$.
The product \eqref{eq:U-count} fixes this normalization independently of the form in which an asymptotic comment is tabulated.

All higher parity-dependent corrections follow from \eqref{eq:finite-dm} with the coefficients \eqref{eq:U-B}. Integer thresholds are obtained from \eqref{eq:residue-rounding} with $r=2$, followed by the minimum over the two parities.

\section{Galois numbers and root-lattice theta sectors}
\label{sec:galois}

\subsection{From one Grassmannian to all subspaces}

The ordinary Galois number is the number of all subspaces of $\F_q^n$:
\begin{equation}\label{eq:Galois-def}
 G_n(q)=\sum_{k=0}^n\qbin{n}{k}.
\end{equation}
At $q=2$ this is A006116 \cite{oeisGalois}:
\[
 1,\ 2,\ 5,\ 16,\ 67,\ 374,\ 2825,\ 29212,\ldots.
\]
The terms near $k=n/2$ have the largest exponential size, but there is no growing Gaussian width: the ratio of the dominant quadratic factors at displacement $j$ is $q^{-j^2}$. Summing these displacements therefore gives a theta constant. The parity of $n$ determines whether the displacements are integral or half-integral.

Define
\begin{equation}\label{eq:theta2}
 \T_0(Q)=\sum_{j\in\Z}Q^{j^2},\qquad
 \T_1(Q)=\sum_{j\in\Z+1/2}Q^{j^2}.
\end{equation}
The leading form is
\begin{equation}\label{eq:Galois-leading}
 G_n(q)\sim\frac{\T_{n\bmod2}(Q)}{\Pinf}\,q^{n^2/4}
 \qquad(n\to\infty\text{ within a fixed parity}).
\end{equation}
Theta leading constants in this enumeration are known \cite{kousidis}. We now prove a stronger exact representation with explicit coefficients at every exponential weight, and extend it at once to flags of arbitrary fixed length.

\subsection{Generalized Galois numbers}

For $r\ge2$, define
\begin{equation}\label{eq:generalGalois-def}
 G_n^{(r)}(q)=\sum_{\substack{k_1,\ldots,k_r\ge0\\k_1+\cdots+k_r=n}}
 \genfrac{[}{]}{0pt}{}{n}{k_1,\ldots,k_r}_q.
\end{equation}
This counts weak flags
\[
 0\subseteq V_1\subseteq\cdots\subseteq V_{r-1}\subseteq\F_q^n;
\]
the successive quotient dimensions are the $k_i$, so repeated subspaces are permitted. For $r=2$ it reduces to \eqref{eq:Galois-def}. For $r=3$, $q=2$, its first terms are
\[
 1,\ 3,\ 12,\ 66,\ 513,\ 5769,\ 95706,\ 2379348,\ldots.
\]
The parameter $r$ is fixed in all asymptotic statements below.

For an integer $\eps$ define the lattice coset
\begin{equation}\label{eq:Lattice}
 \mathcal L_{\eps}^{(r)}=
 \left\{(z_1-\eps/r,\ldots,z_r-\eps/r):
 z_i\in\Z,\ \sum_i z_i=\eps\right\}.
\end{equation}
It lies in the hyperplane $\sum_i j_i=0$. Increasing $\eps$ by $r$ leaves the set unchanged, so the subscript is understood modulo $r$. For $\eps=0$, this is the $A_{r-1}$ root lattice in its usual Euclidean realization. Put
\begin{equation}\label{eq:root-theta}
 \T_{\eps}^{(r)}(Q)=
 \sum_{\boldsymbol j\in\mathcal L_{\eps}^{(r)}}
 Q^{\|\boldsymbol j\|^2/2}>0.
\end{equation}
These series converge absolutely, since a positive definite quadratic form dominates the number of lattice points in successive shells. Reflection of the lattice gives
$\T_{\eps}^{(r)}=\T_{r-\eps}^{(r)}$. At $r=2$, definition \eqref{eq:root-theta} is exactly \eqref{eq:theta2}.

\subsection{An entire function containing all correction coefficients}

Let
\begin{equation}\label{eq:comp-coeff}
 A_{r,M}(Q)=
 \sum_{\substack{m_1,\ldots,m_r\ge0\\m_1+\cdots+m_r=M}}
 \prod_{i=1}^r\frac1{\pp{m_i}},\qquad A_{r,0}=1.
\end{equation}
This is a finite sum. Define
\begin{equation}\label{eq:h-coeff}
 h_{\eps,M}^{(r)}=
 (-1)^M Q^{M^2/(2r)+M/2}\T_{\eps+M}^{(r)}(Q)\,A_{r,M}(Q),
\end{equation}
and
\begin{equation}\label{eq:H-entire}
 H_{\eps}^{(r)}(t)=\sum_{M\ge0}h_{\eps,M}^{(r)}t^M.
\end{equation}

\begin{theorem}[Exact theta-tail representation]\label{thm:theta-tail}
The function $H_{\eps}^{(r)}$ is entire. For every nonnegative integer $n\equiv\eps\pmod r$,
\begin{equation}\label{eq:Galois-exact}
 \boxed{\displaystyle
 G_n^{(r)}(q)=\Pinf^{1-r}
 q^{(r-1)n^2/(2r)}
 \exp\!\bigl(\A(t^r)\bigr)H_{\eps}^{(r)}(t),
 \qquad t=Q^{n/r}.}
\end{equation}
An equivalent entire representation of the tail is
\begin{equation}\label{eq:H-bilateral}
 H_{\eps}^{(r)}(t)=
 \sum_{\boldsymbol j\in\mathcal L_{\eps}^{(r)}}
 Q^{\|\boldsymbol j\|^2/2}
 \prod_{i=1}^r(QtQ^{j_i};Q)_\infty.
\end{equation}
The series in \eqref{eq:H-bilateral} converges absolutely and locally uniformly for all complex $t$.
\end{theorem}

\begin{proof}
Start with the multinomial product
\[
 \genfrac{[}{]}{0pt}{}{n}{k_1,\ldots,k_r}_q
 =q^{(n^2-\sum_i k_i^2)/2}
     \frac{\pp{n}}{\prod_i\pp{k_i}}.
\]
Write $j_i=k_i-n/r$. Then $\boldsymbol j\in\mathcal L_{\eps}^{(r)}$ and
\begin{equation}\label{eq:quadratic-completion}
 \frac{n^2-\sum_i k_i^2}{2}
 =\frac{r-1}{2r}n^2-\frac12\|\boldsymbol j\|^2.
\end{equation}
Using \eqref{eq:product-interpolation}, each denominator contributes
$\Pinf^{-1}(Q^{k_i+1};Q)_\infty$, and the numerator contributes
$\Pinf\exp\A(Q^n)$. This gives \eqref{eq:Galois-exact} with the sum in \eqref{eq:H-bilateral} initially restricted to $k_i\ge0$.

At $t=Q^{n/r}$, a lattice point with any $k_i<0$ contributes zero: the product $(Q^{k_i+1};Q)_\infty$ contains the factor $1-Q^0$. Thus the restricted sum may be extended to the whole lattice coset, once convergence is established.

Expand each product in \eqref{eq:H-bilateral} by \eqref{eq:euler-entire}. For a multi-index $\boldsymbol m=(m_1,\ldots,m_r)$ with $M=\sum_i m_i$, its contribution contains the exponent
\[
 \frac12\|\boldsymbol j\|^2
 +\sum_i m_i j_i+\frac12\sum_i m_i(m_i+1).
\]
Complete the square:
\begin{equation}\label{eq:theta-complete}
 \frac12\left\|\boldsymbol j+\boldsymbol m-\frac Mr\boldsymbol1\right\|^2
 +\frac{M^2}{2r}+\frac M2.
\end{equation}
Translation by $\boldsymbol m-(M/r)\boldsymbol1$ carries
$\mathcal L_{\eps}^{(r)}$ bijectively to $\mathcal L_{\eps+M}^{(r)}$.
The lattice sum at fixed $\boldsymbol m$ is therefore exactly
$Q^{M^2/(2r)+M/2}\T_{\eps+M}^{(r)}$.
Summing over compositions of $M$ proves the coefficient formula \eqref{eq:h-coeff}.

Here is a simultaneous justification of all rearrangements. Since
$\pp{m_i}\ge\Pinf$, and there are $\binom{M+r-1}{r-1}$ weak compositions of $M$,
\begin{equation}\label{eq:H-bound}
 |h_{\eps,M}^{(r)}|
 \le\frac{\max_{0\le j<r}\T_j^{(r)}}{\Pinf^r}
 \binom{M+r-1}{r-1}Q^{M^2/(2r)+M/2}.
\end{equation}
The sum of the right-hand side times $R^M$ converges for every finite $R$. Applying the same completion of squares to the absolute values of the expanded terms proves absolute convergence of the full multiple series, uniformly for $|t|\le R$. Tonelli's theorem and uniform convergence now validate the expansions, the rearrangements, and the entire character of $H$. This also justifies the zero-extension argument at integer $n$.
\end{proof}

\subsection{Forward coefficients, logarithmic coefficients, and the first terms}

Normalize by the leading theta constant. From \eqref{eq:Galois-exact},
\begin{equation}\label{eq:Galois-forward-normalized}
 G_n^{(r)}(q)=\frac{\T_{\eps}^{(r)}}{\Pinf^{r-1}}
 q^{(r-1)n^2/(2r)}\sum_{m\ge0}v_{\eps,m}^{(r)}Q^{mn/r},
\end{equation}
where the series converges and
\begin{equation}\label{eq:Galois-v}
 v_{\eps,m}^{(r)}=
 \frac1{\T_{\eps}^{(r)}}
 \sum_{j=0}^{\lfloor m/r\rfloor}
 \frac{Q^j}{\pp{j}}\,h_{\eps,m-rj}^{(r)}.
\end{equation}
This is a finite formula using only $r$ theta constants and finite composition sums.

For inversion, define
\begin{equation}\label{eq:Galois-B}
 B_{\eps}^{(r)}(t)=\A(t^r)+
       \log\frac{H_{\eps}^{(r)}(t)}{\T_{\eps}^{(r)}}.
\end{equation}
Since $H(0)=\T_{\eps}^{(r)}>0$, the logarithm is analytic in a nonzero disk about zero. If
$c_m=h_{\eps,m}^{(r)}/\T_{\eps}^{(r)}$ for $m\ge1$, then its coefficients are
\begin{equation}\label{eq:Galois-b}
 b_m=\one_{r\mid m}\alpha_{m/r}
 +\sum_{k=1}^m\frac{(-1)^{k+1}}k
  \sum_{\substack{j_1,\ldots,j_k\ge1\\\sum j_i=m}}
       c_{j_1}\cdots c_{j_k}.
\end{equation}
Thus the entire-function representation is not hiding any recurrence in the logarithmic or inverse coefficients.

The first two normalized entire coefficients are
\begin{align}
 c_1&=-\frac{rQ^{(r+1)/(2r)}}{1-Q}
               \frac{\T_{\eps+1}^{(r)}}{\T_{\eps}^{(r)}},
 \label{eq:Galois-c1}\\
 c_2&=Q^{1+2/r}\frac{\T_{\eps+2}^{(r)}}{\T_{\eps}^{(r)}}
 \left(\frac r{(1-Q)(1-Q^2)}
       +\frac{\binom r2}{(1-Q)^2}\right).
 \label{eq:Galois-c2}
\end{align}
Consequently
\begin{equation}\label{eq:Galois-b12}
 b_1=c_1,\qquad b_2=c_2-\frac{c_1^2}2+\one_{r=2}\alpha_1.
\end{equation}
The special contribution at $r=2$ comes from the numerator product $\exp\A(t^r)$.

For ordinary Galois numbers, \eqref{eq:h-coeff} becomes the particularly compact formula
\begin{equation}\label{eq:H-r2}
 h_{\eps,M}^{(2)}=(-1)^M Q^{M(M+2)/4}\T_{\eps+M}(Q)
     \sum_{j=0}^M\frac1{\pp{j}\pp{M-j}}.
\end{equation}
All higher exponentially small corrections to \eqref{eq:Galois-leading} are therefore explicit.

\subsection{Inverse theta transseries}

Define the residue-sheet interpolation by replacing $n$ by a real $x$ on the right-hand side of \eqref{eq:Galois-exact}, with $\eps$ held fixed. It is positive and increasing for all sufficiently large $x$. Its phase data are
\begin{equation}\label{eq:Galois-phase}
 a=\frac{(r-1)h}{2r},\quad b=0,\quad
 c=\log\frac{\T_{\eps}^{(r)}}{\Pinf^{r-1}},\quad
 \lambda=\frac hr.
\end{equation}
In particular,
\begin{equation}\label{eq:Galois-center}
 s_{\eps}=\sqrt{\frac{2r}{(r-1)h}
       \log\left(\frac{Y\Pinf^{r-1}}{\T_{\eps}^{(r)}}\right)},
 \qquad D_{\eps}=\frac{(r-1)h}{r}s_{\eps}.
\end{equation}
The complete inverse is
\begin{equation}\label{eq:Galois-inverse-full}
 \boxed{\displaystyle
 (G_{q,\eps}^{(r)})^{-1}(Y)
 =s_{\eps}+\sum_{m\ge1}d_{\eps,m}(D_{\eps})Q^{ms_{\eps}/r},}
\end{equation}
with \eqref{eq:Galois-b} substituted into \eqref{eq:finite-dm}. The expansion converges for all sufficiently large $Y$.

Its first term is
\begin{align}\label{eq:Galois-first}
 (G_{q,\eps}^{(r)})^{-1}(Y)=s_{\eps}
 &+\frac{rQ^{(r+1)/(2r)}}{(1-Q)D_{\eps}}
     \frac{\T_{\eps+1}^{(r)}}{\T_{\eps}^{(r)}}
      Q^{s_{\eps}/r}
 +O(Q^{2s_{\eps}/r}/D_{\eps}).
\end{align}
This correction is positive. The next coefficient is obtained from \eqref{eq:d2} and \eqref{eq:Galois-b12}, and every further coefficient is still a finite expression.

At $Q=1/2$ the constants are numerically
\begin{align*}
 \T_0^{(2)}&=2.1289368272118771586694585485\ldots,\\
 \T_1^{(2)}&=2.1289312505130275585916134026\ldots,\\
 \T_0^{(3)}&=5.2335188744705752675068634419\ldots,\\
 \T_1^{(3)}=\T_2^{(3)}&=5.2335186066094946295319364209\ldots.
\end{align*}
Their small numerical differences must not be discarded in an all-order inverse calculation. For fixed $Q$, a nonzero change of $c$ shifts the center by order $1/s$, which eventually dominates every exponential sector.

\subsection{Weighted flags and Rogers--Szeg\H{o} polynomials}
\label{sec:weighted-Galois}

A multivariate extension is almost automatic. Fix positive real weights $z_1,\ldots,z_r$ and define
\begin{equation}\label{eq:weightedGalois}
 G_n^{(r)}(q;\boldsymbol z)=
 \sum_{k_1+\cdots+k_r=n}
 \genfrac{[}{]}{0pt}{}{n}{k_1,\ldots,k_r}_q
 \prod_i z_i^{k_i}.
\end{equation}
Replace the theta constant by
\begin{equation}\label{eq:weightedtheta}
 \T_{\eps}^{(r)}(Q;\boldsymbol z)=
 \sum_{\boldsymbol j\in\mathcal L_{\eps}^{(r)}}
 Q^{\|\boldsymbol j\|^2/2}\prod_i z_i^{j_i}.
\end{equation}
The Gaussian factor guarantees convergence for every fixed positive weight vector. The proof of Theorem~\ref{thm:theta-tail}, including its absolute-convergence argument, gives the entire coefficients
\begin{align}\label{eq:weightedh}
 h_{\eps,M}^{(r)}(\boldsymbol z)
 ={}&(-1)^M Q^{M^2/(2r)+M/2}
       \T_{\eps+M}^{(r)}(Q;\boldsymbol z)\\[-0.3em]
 &\times\sum_{m_1+\cdots+m_r=M}
       \prod_{i=1}^r\frac{z_i^{M/r-m_i}}{\pp{m_i}}.\notag
\end{align}
The extra factor is exactly what results from translating the linear weight under the completion of squares \eqref{eq:theta-complete}. The leading prefactor acquires
$\prod_i z_i^{n/r}$, and the inverse phase has
\begin{equation}\label{eq:weightedphase}
 a=\frac{(r-1)h}{2r},\qquad
 b=\frac1r\sum_i\log z_i,\qquad
 c=\log\frac{\T_{\eps}^{(r)}(Q;\boldsymbol z)}{\Pinf^{r-1}}.
\end{equation}
Nothing else changes in the inversion theorem.

At $r=2$, $z_1=z$, $z_2=1$, the finite sum is the Rogers--Szeg\H{o} polynomial
$\sum_k\qbin{n}{k}z^k$. The relation between generalized Galois numbers, flags, and generalized Rogers--Szeg\H{o} polynomials is also studied in \cite{bliem}. Here fixed positive weights preserve an explicit convergent inverse transseries with the same small parameter $Q^{s/r}$.

\section{Inverting convergence to a finite limit}
\label{sec:endpoint}

The previous examples grow without bound, but their normalized versions often converge to positive constants. Their inverse problem asks how large an index is needed to reach a small prescribed deficit. The correct leading inverse is logarithmic.

\subsection{A general endpoint inversion formula}

\begin{theorem}[Logarithmic endpoint inverse]\label{thm:endpoint}
Let $B(t)=b_1t+b_2t^2+\cdots$ be analytic at zero with $b_1>0$, and set
$\eps=B(\ee^{-\lambda x})$. For positive sufficiently small $\eps$, the large real solution has a convergent expansion
\begin{equation}\label{eq:endpoint-all}
 \boxed{\displaystyle
 x=\frac1\lambda\log\frac{b_1}{\eps}
 -\frac1\lambda\sum_{m\ge1}\frac{\eps^m}{m}
     \coef{u}{m}\left(\frac{u}{B(u)}\right)^m.}
\end{equation}
In particular,
\begin{equation}\label{eq:endpoint-first}
 x=\frac1\lambda\log\frac{b_1}{\eps}
  +\frac{b_2}{\lambda b_1^2}\eps
  +\frac1\lambda\left(\frac{b_3}{b_1^3}
             -\frac{3b_2^2}{2b_1^4}\right)\eps^2
  +O(\eps^3).
\end{equation}
\end{theorem}
\begin{proof}
The analytic inverse theorem gives $t=t(\eps)$ with $B(t)=\eps$ and
$t=\eps/b_1+O(\eps^2)$. Let $R(t)=t/B(t)$, analytic at zero. Since
$t/\eps=R(t)$, we have
$x=-\lambda^{-1}\log\eps-\lambda^{-1}\log R(t(\eps))$.
Lagrange--B\"urmann gives, for $m\ge1$,
\[
 [\eps^m]\log R(t(\eps))
 =\frac1m[u^{m-1}]\frac{R'(u)}{R(u)}R(u)^m
 =\frac1m[u^m]R(u)^m.
\]
The constant term is $\log R(0)=-\log b_1$. This proves \eqref{eq:endpoint-all}; expanding the first two coefficients proves \eqref{eq:endpoint-first}. All series used are analytic near zero, establishing convergence.
\end{proof}

\subsection{Invertibility probabilities of random matrices}

The probability that a uniformly random $n\times n$ matrix over $\F_q$ is invertible is
\begin{equation}\label{eq:probability}
 p_n(q)=\frac{|\GL(n,q)|}{q^{n^2}}=\pp{n}.
\end{equation}
It decreases to $\Pinf$. With the canonical interpolation and the logarithmic relative excess
\begin{equation}\label{eq:probability-eps}
 \eps=\log\frac{\mathcal P_Q(x)}{\Pinf},
\end{equation}
we have exactly $\eps=\A(Q^x)$. Therefore
\begin{equation}\label{eq:probability-inverse}
 x=\frac1h\log\frac{\alpha_1}{\eps}
 +\frac{\alpha_2}{h\alpha_1^2}\eps
 +\frac1h\left(\frac{\alpha_3}{\alpha_1^3}
       -\frac{3\alpha_2^2}{2\alpha_1^4}\right)\eps^2
 +O(\eps^3).
\end{equation}
At $q=2$, this becomes
\begin{equation}\label{eq:probability-q2}
 x=\frac1{\log2}\left(
       \log\frac1\eps+\frac\eps6+\frac{\eps^2}{168}
       +O(\eps^3)\right).
\end{equation}
For a relative excess $\Delta=\mathcal P_Q(x)/\Pinf-1$, substitute
$\eps=\log(1+\Delta)$, not $\eps=\Delta$. Confusing these two small variables changes every coefficient after the leading approximation.

Since the interpolation decreases, the least integer $n$ satisfying
$p_n(q)\le\Pinf\ee^{\eps}$ is the ceiling of this real inverse, provided it lies in the nonnegative range. Strict inequalities require the corresponding strict-rounding convention.

\subsection{Normalized central Gaussian coefficients}

Define
\begin{equation}\label{eq:central-normalized}
 R_q(x)=\Pinf q^{-x^2}H_{1,1}(x)
       =\exp\bigl(\A(t^2)-2\A(t)\bigr),\qquad t=Q^x.
\end{equation}
For sufficiently large $x$, this approaches one from below. Indeed,
\[
 -\log R_q(x)=2\A(t)-\A(t^2)=2\alpha_1t+O(t^2)>0.
\]
Its derivative with respect to $t$ is positive near zero, so $R_q$ increases to one there. Theorem~\ref{thm:endpoint} applies to
$\eps=-\log R_q(x)$ with
\begin{equation}\label{eq:central-endpoint-b}
 b_m=2\alpha_m-\one_{2\mid m}\alpha_{m/2}.
\end{equation}
It gives an explicit logarithmic inverse for attaining a prescribed normalized accuracy, independently of the growing-count inverse in Section~\ref{sec:gaussian}.

\subsection{Higher-order leading tails and Puiseux corrections}

Suppose instead
\[
 B(t)=b_rt^r(1+O(t)),\qquad b_r>0.
\]
Set $\eta=(\eps/b_r)^{1/r}$ and choose the analytic root
\[
 C(t)=\left(B(t)/b_r\right)^{1/r}=t(1+O(t)).
\]
The ordinary analytic inverse of $C$ gives $t=\eta(1+O(\eta))$. Hence
\begin{equation}\label{eq:puiseux-endpoint}
 x=\frac1{r\lambda}\log\frac{b_r}{\eps}
       +\sum_{m\ge1}\gamma_m\eps^{m/r},
\end{equation}
with coefficients obtained by the same Lagrange formula applied to $C$. The series in $\eps^{1/r}$ converges near zero on the chosen positive branch.

There is an important simplification when $B(t)=\widehat B(t^r)$: put $u=t^r$ and invert $\eps=\widehat B(u)$ directly. Only integer powers of $\eps$ remain after the logarithm. In particular the normalized symplectic count has this simpler structure, because its tail is $\mathcal A_{Q^2}(Q^{2x})$. Its leading index is
\[
 x\sim\frac1{2h}\log\frac{Q^2}{(1-Q^2)\eps},
\]
followed by a convergent series in integer powers of $\eps$.

\section{Primitive necklaces and irreducible polynomials}
\label{sec:arithmetic}

\subsection{Exact divisor sectors}

Let $I_q(n)$ be the number of monic irreducible polynomials of degree $n$ over $\F_q$, for prime-power $q$. For any integer alphabet size $q\ge2$, the same formula counts primitive necklaces, or equivalently Lyndon words, of length $n$:
\begin{equation}\label{eq:irreducibles}
 I_q(n)=\frac1n\sum_{d\mid n}\mu(d)q^{n/d},\qquad n\ge1.
\end{equation}
The binary specialization is A001037 \cite{oeisI}. Its conventional entry at $n=0$ is not used here.

Here are two proofs of the formula. The polynomial $X^{q^n}-X$ is the product of all monic irreducibles whose degrees divide $n$, with no repeated roots. Taking degrees gives
$q^n=\sum_{d\mid n}dI_q(d)$, and M\"obius inversion yields \eqref{eq:irreducibles}. For words, every word has a unique least period $d\mid n$. A primitive rotation orbit of length $d$ contains exactly $d$ words, giving the same identity and the same inversion formula.

Normalize by the leading term:
\begin{equation}\label{eq:arithmetic-sectors}
 I_q(n)=\frac{q^n}{n}
 \left(1+\sum_{d\ge2}\mu(d)\one_{d\mid n}
       \ee^{-h(1-1/d)n}\right).
\end{equation}
At each integer $n$ the sum is finite. The exponential actions are
$h(1-1/d)$, with coefficients modulated by divisibility. They accumulate at $h$ as $d\to\infty$.

For a fixed cutoff $D_0\ge2$, the omitted relative tail satisfies
\begin{equation}\label{eq:arithmetic-cutoff}
 \left|\sum_{\substack{d\mid n\\d>D_0}}
       \mu(d)\ee^{-h(1-1/d)n}\right|
 \le n\ee^{-h(1-1/(D_0+1))n}.
\end{equation}
Indeed, there are at most $n$ possible divisors and each exponent is at least the displayed action. This gives a rigorous finite-sector asymptotic decomposition uniformly over all integers $n$.

It is not, however, a single analytic power series in $\ee^{-\lambda n}$ with fixed constant coefficients. Along even $n$ the first relative correction is
\[
 -q^{-n/2}+O(nq^{-2n/3}),
\]
whereas along prime $n$ the exact correction is $-q^{1-n}$. The lowest nonzero action changes with the arithmetic subsequence. On a fixed congruence class modulo $\operatorname{lcm}(1,\ldots,D_0)$, the divisibility coefficients up to $D_0$ are fixed, but this does not fix the later coefficients.

There is a distinction here between an arithmetic sector decomposition and a particular formal transseries field. The accumulation of actions does not by itself exclude every possible Hahn-series construction. The concrete point is that no fixed finite exponential grid and no fixed analytic coefficient germ describes all the divisibility data in \eqref{eq:arithmetic-sectors}.

\subsection{A finite-radical interpolation removes the moving divisors}

There is nevertheless a natural class of subsequences on which a full analytic inverse transseries can be constructed. Fix a squarefree integer $R\ge2$. Whenever $\rad(n)=R$, all nonzero M\"obius terms in \eqref{eq:irreducibles} correspond to divisors of $R$. Define
\begin{equation}\label{eq:ER}
 E_R(t)=\sum_{d\mid R}\mu(d)t^{R-R/d}
       =1+\sum_{\substack{d\mid R\\d>1}}\mu(d)t^{R-R/d},
 \qquad \kappa=\frac hR.
\end{equation}
This is a finite polynomial with constant term one. The analytic function
\begin{equation}\label{eq:fixedradical-function}
 J_{q,R}(x)=\frac{\ee^{hx}}x E_R(\ee^{-\kappa x})
          =\frac1x\sum_{d\mid R}\mu(d)\ee^{hx/d}
\end{equation}
agrees with $I_q(n)$ at every $n$ with radical $R$. It is positive and increasing for sufficiently large $x$.

The function is now explicitly specified, so its inverse has a well-defined exponentially small structure. We use
\begin{equation}\label{eq:logphase}
 P(x)=hx-\log x,\qquad B(t)=\log E_R(t),\qquad
 D=P'(s)=h-\frac1s.
\end{equation}
The large real center is
\begin{equation}\label{eq:Wcenter}
 \boxed{\displaystyle
 s=-\frac1h W_{-1}\!\left(-\frac hY\right).}
\end{equation}
The $W_{-1}$ branch is required because $s\to\infty$. The principal branch would give the small solution of $\ee^{hx}/x=Y$. The branch convention agrees with DLMF and \texttt{ProductLog[-1,z]} \cite{dlmfW,wolframW}.

\subsection{A general inverse formula beyond quadratic phases}

The coefficient extraction mechanism does not depend on $P$ being quadratic. Let $X(v)$ be the local inverse of $P$ at $v=P(s)$, and define
\begin{equation}\label{eq:generalPhi}
 \Phi_s(u)=X(P(s)-B(u))-s.
\end{equation}
Exactly as in Section~\ref{sec:engine}, putting
$u=t\ee^{-\kappa\delta}$ gives $\delta=\Phi_s(u)$. Hence
\begin{equation}\label{eq:general-inverse-extraction}
 d_m(s)=-\frac1{m\kappa}[u^m]\exp\bigl(-m\kappa\Phi_s(u)\bigr).
\end{equation}
For \eqref{eq:logphase}, this is explicit in Lambert $W$:
\begin{equation}\label{eq:W-Phi}
 \Phi_s(u)=-\frac1h W_{-1}\!\left(-\frac hY E_R(u)\right)-s,
\end{equation}
where the branch is the local analytic continuation of its real value at $u=0$.

There is also a finite derivative formula that avoids expanding a special function. Let $b_m=[u^m]B(u)$ and let $C_{m,k}$ be as in \eqref{eq:Cmk}. Define an operator on functions of $s$ by
\begin{equation}\label{eq:shiftedoperator}
 \mathcal D_m=\frac1{P'(s)}\left(\frac{\dd}{\dd s}-m\kappa\right).
\end{equation}
Then
\begin{equation}\label{eq:general-derivative}
 \boxed{\displaystyle
 d_m(s)=\sum_{k=1}^m\frac{(-1)^kC_{m,k}}{k!}
             \mathcal D_m^{\,k-1}\left(\frac1{P'(s)}\right).}
\end{equation}
Only finitely many differentiations are involved at a prescribed weight.

\begin{proof}[Proof of the derivative formula]
Temporarily insert a scalar parameter $\eta$ in
$P(x)+\eta B(z\ee^{-\kappa x})=P(s)$, and put $v=P(x)$. Ordinary Lagrange inversion for
$v+\eta\widetilde B(v)=P(s)$, followed by the analytic function $x=X(v)$, gives the coefficient of $\eta^k$ as
\[
 \frac{(-1)^k}{k!}
 \left(\frac{\dd}{\dd v}\right)^{k-1}
 \left(X'(v)\widetilde B(v)^k\right)_{v=P(s)}.
\]
Extract the coefficient of $z^m$. It is
\[
 \frac{(-1)^k C_{m,k}}{k!}
 \left(\frac1{P'(s)}\frac{\dd}{\dd s}\right)^{k-1}
 \left(\frac{\ee^{-m\kappa s}}{P'(s)}\right).
\]
Conjugating the differential operator by $\ee^{-m\kappa s}$ gives \eqref{eq:shiftedoperator}. For fixed $m$, terms with $k>m$ cannot contribute. Set $\eta=1$ and remove the common factor $\ee^{-m\kappa s}$ to obtain \eqref{eq:general-derivative}.
\end{proof}

\begin{theorem}[Inverse transseries on a fixed-radical sheet]\label{thm:radical-inverse}
For fixed $q>1$ and squarefree $R\ge2$, the large real inverse of \eqref{eq:fixedradical-function} has a convergent expansion
\begin{equation}\label{eq:radical-inverse}
 J_{q,R}^{-1}(Y)=s+\sum_{m\ge1}d_m(s)\ee^{-mhs/R},
\end{equation}
where $s$ is \eqref{eq:Wcenter} and the coefficients are given by
\eqref{eq:general-inverse-extraction} or \eqref{eq:general-derivative}. Each coefficient is a rational function of $s,h,\kappa$ and the finitely many coefficients of $B$ needed at that weight.
\end{theorem}
\begin{proof}
For $s\ge2/h$, $P'(s)\ge h/2$. On a fixed disk $|\delta|\le\rho$ with $s$ large,
\[
 P(s+\delta)-P(s)-P'(s)\delta=O(\delta^2/s^2),
\]
uniformly in $s$. This follows by expanding $-\log(1+\delta/s)$, or from its second derivative. Since $B$ is analytic near zero, the same Rouch\'e argument used in Theorem~\ref{thm:inverse} gives a holomorphic solution in a fixed disk of the auxiliary small parameter $t$, uniformly for sufficiently large $s$. Substituting $t=\ee^{-\kappa s}$ proves convergence. Formula \eqref{eq:general-derivative}, together with $P'(s)=h-1/s$, proves rationality of the coefficients. Eventual monotonicity follows from $P'(x)\to h>0$ and the exponentially small tail derivative.
\end{proof}

\subsection{Prime-power indices: an explicit first example}

If $n=p^a$ with a fixed prime $p$, then
\[
 I_q(n)=\frac{q^n-q^{n/p}}n.
\]
It is convenient to use the primitive action
$\kappa=h(1-1/p)$ rather than $h/p$. The selected interpolation is
\begin{equation}\label{eq:primepower-interp}
 J_{q,p}(x)=\frac{\ee^{hx}}x(1-\ee^{-\kappa x}),
 \qquad B(t)=\log(1-t).
\end{equation}
With $D=h-1/s$ and $t_0=\ee^{-\kappa s}$, the first two inverse sectors are
\begin{equation}\label{eq:primepower-inverse}
 \boxed{\displaystyle
 J_{q,p}^{-1}(Y)=s+\frac{t_0}{D}
 +\left(\frac1{2D}-\frac{\kappa}{D^2}
              -\frac1{2s^2D^3}\right)t_0^2+O(t_0^3).}
\end{equation}
To check the second coefficient directly, expand
\[
 P(s+\delta)-P(s)=D\delta+\frac{\delta^2}{2s^2}+O(\delta^3),
 \qquad B(t)=-t-\frac{t^2}{2}+O(t^3),
\]
and substitute into the inverse equation. The coefficient of $t_0$ is $1/D$; the coefficient of $t_0^2$ is exactly the bracket in \eqref{eq:primepower-inverse}.

The exponential parameter itself can be written in terms of $Y$ and the exact Lambert center:
\begin{equation}\label{eq:primepower-scale}
 t_0=(Ys)^{-(1-1/p)},\qquad \ee^{hs}=Ys.
\end{equation}
Thus the inverse combines a logarithmic Lambert-$W$ scale with fractional inverse powers of $Y$ and corresponding powers of $s$. Writing
$\ell=\log(Y/h)$ and $v=\log\ell$, the beginning of the ordinary logarithmic expansion is
\begin{equation}\label{eq:W-logexpansion}
 hs=\ell+v+\frac v\ell+\frac{v-v^2/2}{\ell^2}
       +O(v^3/\ell^3).
\end{equation}
It follows by substitution into $hs-\log(hs)=\ell$; the full logarithmic calculus is already developed in the source article \cite{base}. Retaining $s$ exactly in \eqref{eq:primepower-inverse} keeps all of those logarithmic corrections organized inside a single center.

For a second example, $R=6$ gives
\begin{equation}\label{eq:R6}
 E_6(t)=1-t^3-t^4+t^5,
 \qquad
 J_{q,6}(x)=\frac{q^x-q^{x/2}-q^{x/3}+q^{x/6}}x.
\end{equation}
The inverse weights in $t_0=\ee^{-hs/6}$ start at $3,4,5$ and include their additive combinations. Weights $1$ and $2$ vanish. The first three nonzero terms are
\begin{equation}\label{eq:R6-first}
 J_{q,6}^{-1}(Y)=s+\frac{t_0^3+t_0^4-t_0^5}{D}+O(t_0^6).
\end{equation}
Nonlinear interactions first enter at weight six.

\subsection{A uniform two-candidate theorem for the integer inverse}

The fixed-radical sheets do not define a single interpolation for all $n$. The integer threshold can nevertheless be controlled uniformly without choosing any interpolation of the arithmetic corrections.

\begin{lemma}[Uniform deficit bound]\label{lem:irreducible-bound}
For every integer $q\ge2$ and $n\ge1$,
\begin{equation}\label{eq:uniform-deficit}
 0\le q^n-nI_q(n)
 \le\sum_{j=1}^{\lfloor n/2\rfloor}q^j
 =\frac{q^{\lfloor n/2\rfloor+1}-q}{q-1}.
\end{equation}
In particular,
\begin{equation}\label{eq:I-lower}
 \frac{q^n}{n}\left(1-\frac q{q-1}q^{-n/2}\right)
 \le I_q(n)\le\frac{q^n}{n}.
\end{equation}
\end{lemma}
\begin{proof}
The deficit counts nonprimitive words. Every nonprimitive word has a proper period dividing $n$, hence a period at most $\lfloor n/2\rfloor$. There are at most $q^j$ words obtained by repeating a block of length $j$. Taking the union bound over all $1\le j\le\lfloor n/2\rfloor$ proves \eqref{eq:uniform-deficit}; allowing lengths that do not divide $n$ merely enlarges the upper bound. Dividing by $q^n$ gives \eqref{eq:I-lower}.
\end{proof}

\begin{theorem}[An exact two-candidate threshold rule]\label{thm:two-candidate}
Let $q\ge2$ be an integer, and let
\[
 X(Y)=-\frac1hW_{-1}(-h/Y),\qquad n_0=\lceil X(Y)\rceil.
\]
If $n_0\ge4$, then the integer threshold
$N_q(Y)=\min\{n\ge1:I_q(n)\ge Y\}$ satisfies
\begin{equation}\label{eq:two-candidate}
 \boxed{\displaystyle
 N_q(Y)=
 \begin{cases}
 n_0,&I_q(n_0)\ge Y,\\
 n_0+1,&I_q(n_0)<Y.
 \end{cases}}
\end{equation}
Thus one exact divisor-sum comparison resolves the threshold after computing the Lambert center.
\end{theorem}
\begin{proof}
Put $B_q(n)=q^n/n$. Its values are nondecreasing for positive integers, and strictly increasing from $n\ge2$, because
$B_q(n+1)/B_q(n)=qn/(n+1)$. The choice of $n_0$ and the large increasing branch of $q^x/x$ imply
$I_q(n)\le B_q(n)<Y$ for every $n<n_0$.

It remains to prove $I_q(n_0+1)\ge Y$. By \eqref{eq:I-lower}, for $n\ge4$,
\begin{align*}
 \frac{I_q(n+1)}{B_q(n)}
 &\ge\frac{qn}{n+1}
   \left(1-\frac q{q-1}q^{-(n+1)/2}\right)\\
 &\ge\frac{2n}{n+1}\left(1-2\cdot2^{-(n+1)/2}\right)>1.
\end{align*}
The last expression increases with $n$ and already exceeds one at $n=4$.
Since $Y\le B_q(n_0)$, this proves $I_q(n_0+1)>Y$. No separate monotonicity theorem for $I_q(n)$ was required.
\end{proof}

The extra candidate cannot be discarded. For $Y=q^n/n$ with $n\ge4$, nonprimitive words exist, so $I_q(n)<Y$ and the answer is $n+1$. On the other hand, at $Y=I_q(n)$ with $n$ sufficiently large the answer is $n$. The switch occurs inside a relative strip of width $O(q^{-n/2})$. This is precisely where uncontrolled rounding of a smooth approximation loses information.

\subsection{All necklaces and a related arithmetic spectrum}

The number of all $q$-color necklaces of length $n$ is
\begin{equation}\label{eq:allnecklaces}
 M_q(n)=\frac1n\sum_{d\mid n}\varphi(d)q^{n/d}.
\end{equation}
Burnside's lemma proves the formula: a rotation by $j$ positions fixes
$q^{\gcd(n,j)}$ words, and grouping rotations with the same greatest common divisor gives \eqref{eq:allnecklaces}. The binary case is A000031 \cite{oeisnecklaces}.

Its relative arithmetic sectors have the same actions as \eqref{eq:arithmetic-sectors}, now with coefficients $\varphi(d)\one_{d\mid n}$, all positive. A uniform elementary estimate is
\[
 0\le\frac{nM_q(n)}{q^n}-1
 \le n^2q^{-n/2}.
\]
There are at most $n$ divisor terms, each $\varphi(d)\le n$, and $n/d\le n/2$ for $d>1$. Thus the Lambert center still gives the leading inverse scale, but its correction is on the opposite side from that for primitive necklaces. Unlike the M\"obius sum, fixing the radical alone does not fix all nonzero totient terms; powers of its primes continue to appear among the divisors. This explains why the finite polynomial $E_R$ construction is specific to the primitive-word problem.

\section{The singular transition \texorpdfstring{$q\to1$}{q to 1}}
\label{sec:double}

\subsection{Why a fixed-\texorpdfstring{$q$}{q} transseries cannot simply be specialized}

The coefficients $\alpha_m=Q^m/[m(1-Q^m)]$ grow like $1/(hm^2)$ as $h\to0$. Simultaneously $\Pinf$ tends to zero exponentially in $1/h$. A limit that lets $q\to1$ while $x\to\infty$ therefore reorganizes both the nominal constant and the finite-product tail.

Consider the scaling
\begin{equation}\label{eq:double-scaling}
 Q=\ee^{-h},\qquad h\downarrow0,\qquad x=\frac\tau h,
 \qquad \tau>0.
\end{equation}
The analytic parameter $q$ is continuous here; there is no family of finite fields whose orders tend to one. The finite-field interpretation motivates the formulas, but this limiting problem concerns their real-$q$ special-function extensions.

\subsection{An all-order dilogarithmic expansion of the tail}

Write $z=\ee^{-\tau}$. The exact tail is
\begin{equation}\label{eq:double-Aexact}
 \mathcal A_{\ee^{-h}}(z)
 =\sum_{m\ge1}\frac{\ee^{-m\tau}}{m(\ee^{hm}-1)}.
\end{equation}
Let the Bernoulli numbers be defined by
$v/(\ee^v-1)=\sum_{j\ge0}B_jv^j/j!$ near zero.

\begin{proposition}[Uniform finite-order tail expansion]\label{prop:double-A}
For every fixed integer $K\ge1$ and every $\tau_0>0$, uniformly for $\tau\ge\tau_0$ as $h\downarrow0$,
\begin{align}\label{eq:double-A}
 \mathcal A_{\ee^{-h}}(\ee^{-\tau})
 ={}&\frac{\Li_2(\ee^{-\tau})}{h}
     +\frac12\log(1-\ee^{-\tau})\\[-0.3em]
 &+\sum_{k=1}^K\frac{B_{2k}}{(2k)!}h^{2k-1}
          \Li_{2-2k}(\ee^{-\tau})
 +O_{K,\tau_0}(h^{2K+1}).\notag
\end{align}
The assertion remains valid after any fixed number of derivatives with respect to $\tau$, with a corresponding change of the implied constant.
\end{proposition}
\begin{proof}
The remainder in the local expansion
\[
 \frac1{\ee^v-1}=\frac1v-\frac12+
   \sum_{k=1}^K\frac{B_{2k}}{(2k)!}v^{2k-1}+r_K(v)
\]
satisfies $|r_K(v)|\le C_Kv^{2K+1}$ for all real $v>0$. Near zero this follows from analyticity after removal of $1/v$. On a compact positive interval it follows from continuity, and for large $v$ the displayed polynomial and $1/v$ grow no faster than a constant times $v^{2K+1}$.

Insert $v=hm$ into \eqref{eq:double-Aexact}. The sum of the absolute remainders is at most
\[
 C_Kh^{2K+1}\sum_{m\ge1}\ee^{-m\tau}m^{2K},
\]
which is uniformly bounded by a constant times $h^{2K+1}$ for $\tau\ge\tau_0$. Summing the explicit powers of $m$ gives the polylogarithms in \eqref{eq:double-A}. Differentiation in $\tau$ adds only a fixed power of $m$, leaving the same argument valid.
\end{proof}

\subsection{The modular factor and its separate exponential scale}

The modular transformation of the Euler product gives the exact identity \cite{dlmfqmod}
\begin{equation}\label{eq:modularP}
 \mathcal P_{\ee^{-h}}
 =\sqrt{\frac{2\pi}{h}}
   \exp\left(-\frac{\pi^2}{6h}+\frac h{24}\right)
   \mathcal P_{\widehat Q},
 \qquad \widehat Q=\ee^{-4\pi^2/h}.
\end{equation}
For example, it follows from the Dedekind eta transformation
$\eta(-1/\tau)=\sqrt{-i\tau}\,\eta(\tau)$ at $\tau=ih/(2\pi)$, together with
$\eta(ih/(2\pi))=\ee^{-h/24}\mathcal P_{\ee^{-h}}$.
The exponentially small product left over is explicit:
\begin{equation}\label{eq:modularsectors}
 -\log\mathcal P_{\widehat Q}
 =\sum_{N\ge1}\sigma_{-1}(N)\widehat Q^N,
 \qquad\sigma_{-1}(N)=\sum_{d\mid N}\frac1d.
\end{equation}
To prove \eqref{eq:modularsectors}, expand each logarithm of the Euler product and group the terms with the same product of indices. This series converges for $0<\widehat Q<1$.

The new variable $\widehat Q$ is not a power of $Q^x=\ee^{-\tau}$. The simultaneous limit therefore contains two different mechanisms: a finite-size tail described by \eqref{eq:double-Aexact}, and the modular exponential factor \eqref{eq:modularsectors}. The latter is an exact separated contribution. It is not being asserted to describe the entire beyond-all-orders structure of the finite-size tail's Bernoulli expansion.

\subsection{Central Gaussian coefficients and the crossover phase}

For the central Gaussian interpolation,
\begin{equation}\label{eq:central-double-exact}
 \log H_{1,1}(\tau/h)
 =\frac{\tau^2}{h}-\log\mathcal P_{\ee^{-h}}
 +\mathcal A_{\ee^{-h}}(\ee^{-2\tau})
 -2\mathcal A_{\ee^{-h}}(\ee^{-\tau}).
\end{equation}
Introduce
\begin{align}
 S(\tau)&=\tau^2+\frac{\pi^2}{6}
               +\Li_2(\ee^{-2\tau})-2\Li_2(\ee^{-\tau}),
 \label{eq:Sphase}\\
 J(\tau)&=\frac12\log(1-\ee^{-2\tau})
                  -\log(1-\ee^{-\tau}),\label{eq:Jamp}\\
 E_1(\tau)&=-\frac1{24}
       +\frac{\Li_0(\ee^{-2\tau})-2\Li_0(\ee^{-\tau})}{12},
 \label{eq:E1}\\
 E_k(\tau)&=\frac{B_{2k}}{(2k)!}
       \left(\Li_{2-2k}(\ee^{-2\tau})
                 -2\Li_{2-2k}(\ee^{-\tau})\right),\quad k\ge2.
 \label{eq:Ek}
\end{align}
Combining \eqref{eq:double-A} and \eqref{eq:modularP} gives
\begin{align}\label{eq:central-double-asympt}
 \log H_{1,1}(\tau/h)
 ={}&\frac{S(\tau)}h-\frac12\log\frac{2\pi}{h}+J(\tau)
     +\sum_{k=1}^K h^{2k-1}E_k(\tau)\\[-0.3em]
 &-\log\mathcal P_{\widehat Q}
     +O_{K,\tau_0}(h^{2K+1}).\notag
\end{align}
The modular term is retained exactly, although a finite algebraic remainder is of course much larger than that term. Equation \eqref{eq:central-double-exact} remains available when the finite-size tail itself must be evaluated beyond that algebraic accuracy.

The derivative of the new dominant phase has a useful elementary form:
\begin{equation}\label{eq:Sderivative}
 \boxed{\displaystyle
 S'(\tau)=2\log(1+\ee^{\tau})>0,\qquad
 S''(\tau)=\frac{2\ee^\tau}{1+\ee^\tau}>0.}
\end{equation}
Indeed, differentiating the dilogarithms gives
$2\tau+2\log(1-\ee^{-2\tau})-2\log(1-\ee^{-\tau})$
which reduces to \eqref{eq:Sderivative}. Also $S(0)=0$ by continuity. Thus $S$ is increasing from zero to infinity and has an unambiguous positive inverse.

For orientation, its endpoint expansions are
\begin{align}
 S(\tau)&=2(\log2)\tau+\frac{\tau^2}{2}
                  +\frac{\tau^3}{12}+O(\tau^5),\quad\tau\to0,
 \label{eq:Ssmall}\\
 S(\tau)&=\tau^2+\frac{\pi^2}{6}+O(\ee^{-\tau}),\quad\tau\to\infty.
 \label{eq:Slarge}
\end{align}
They follow from \eqref{eq:Sderivative} and the defining dilogarithms. The small-$\tau$ phase and $J(\tau)\sim\tfrac12\log(2/\tau)$ display the classical central-binomial phase and amplitude. The uniform remainder theorem above is stated for $\tau\ge\tau_0>0$; matching it uniformly all the way to $\tau=0$ requires a different error analysis.

\subsection{Inverting the crossover phase}

Let $\ell=h\log Y$ be fixed and positive, and define
\begin{equation}\label{eq:crossover-center}
 \tau_*=S^{-1}(\ell),\qquad
 L_h=\frac12\log\frac{2\pi}{h}.
\end{equation}
On compact positive ranges of $\ell$, the derivative in \eqref{eq:Sderivative} is bounded away from zero. The inverse index therefore has the expansion
\begin{equation}\label{eq:crossover-inverse}
 x=\frac{\tau_*}{h}+c_1(\tau_*,L_h)
                         +h c_2(\tau_*,L_h)
                         +O\bigl(h^2(1+L_h)^3\bigr),
\end{equation}
where
\begin{align}
 c_1&=\frac{L_h-J(\tau_*)}{S'(\tau_*)},\label{eq:crossover-c1}\\
 c_2&=-\frac{\tfrac12 S''(\tau_*)c_1^2
                  +J'(\tau_*)c_1+E_1(\tau_*)}{S'(\tau_*)}.
 \label{eq:crossover-c2}
\end{align}
To prove these formulas, write $\tau=\tau_*+hc_1+h^2c_2+\cdots$, multiply \eqref{eq:central-double-asympt} by $h$, and compare coefficients in
\[
 S(\tau)+h(J(\tau)-L_h)+h^2E_1(\tau)+\cdots=\ell.
\]
The error estimate follows from Proposition~\ref{prop:double-A}, Taylor's theorem, and the positive derivative. The modular term is smaller than every fixed power of $h$ and so does not affect this finite algebraic estimate.

There is a finite coefficient formula at every higher order. Treat $L_h$ as an independent symbol during coefficient extraction, and put
\begin{equation}\label{eq:crossoverR}
 \mathcal R(h,\tau)=h(J(\tau)-L_h)
                  +\sum_{j\ge1}h^{2j}E_j(\tau).
\end{equation}
The formal series $\tau=\tau_*+\sum_{m\ge1}c_mh^m$ has
\begin{equation}\label{eq:crossover-all}
 c_m=\sum_{k=1}^m\frac{(-1)^k}{k!}
 \left[
 \left(\frac1{S'(\tau)}\frac{\dd}{\dd\tau}\right)^{k-1}
 \frac{[h^m]\mathcal R(h,\tau)^k}{S'(\tau)}
 \right]_{\tau=\tau_*}.
\end{equation}
The proof is the same inverse-perturbation identity used for
\eqref{eq:general-derivative}. At each fixed $m$, only finitely many $E_j$ occur. Each $c_m$ is a polynomial in $L_h$ of degree at most $m$, with smooth coefficients in $\tau_*$. The finite-order expansions are asymptotic with remainders controlled by the preceding propositions; no convergence assertion is made for this Bernoulli-based expansion in $h$.

The main lesson is quantitative. At fixed positive $\ell=h\log Y$, the correct leading center is $S^{-1}(\ell)/h$, not the fixed-$q$ square-root center with its finite-size tail dropped. The tail contributes at order $1/h$ in this regime and therefore belongs to the dominant phase itself.

\section{Error certification, threshold decisions, and computation}
\label{sec:certification}

\subsection{Analytic error bounds versus numerical evidence}

Theorem~\ref{thm:inverse} supplies an analytic, geometrically decreasing outer-sector remainder. For a concrete parameter value, it is often more economical to certify an approximation by its residual. Let
\begin{equation}\label{eq:residual}
 \mathcal E(x)=P(x)+B(\ee^{-\lambda x})-\log Y.
\end{equation}
Suppose a known interval $[u,v]$ contains the desired root and
$\mathcal E'(x)\ge\mu>0$ throughout it. Then for any $x_M\in[u,v]$,
\begin{equation}\label{eq:residualbound}
 |x_M-F^{-1}(Y)|\le\frac{|\mathcal E(x_M)|}{\mu}.
\end{equation}
This is the mean value theorem. Alternatively, checking opposite signs at
$x_M-\eta$ and $x_M+\eta$, together with monotonicity, directly certifies a root enclosure of radius $\eta$.

For $q$-product tails, Lemma~\ref{lem:Atail} bounds the omitted logarithmic terms using only elementary quantities. Derivatives can be bounded in the same way after differentiating the absolutely convergent series. For a theta-tail function, the majorant \eqref{eq:H-bound} bounds the omitted entire coefficients; a fixed polynomial growth factor is dominated by the Gaussian decrease in their weight.

The supplied program uses high-precision floating-point arithmetic for its analytic tests. Those tests are reproducibility checks, not interval-arithmetic certificates. The proofs and inequalities above establish the analytical assertions independently of floating-point evidence. For a certified implementation one would evaluate the residual, derivative bounds, and product or theta remainders with outward-rounded interval arithmetic.

\subsection{A safe ceiling criterion}

If a real inverse has a certified enclosure $[u,v]$ and
$\lceil u\rceil=\lceil v\rceil$, its integer threshold is determined. If the enclosure crosses an integer, evaluate the original exact sequence at that candidate integer. The same rule applies to a residue-class lattice, using \eqref{eq:residue-rounding}.

It is essential to retain equality cases. For a threshold defined using $a_n\ge Y$, the correct rounding is a ceiling of the real inverse. For a condition $a_n>Y$, an exact integer inverse instead requires the next lattice point. A floating-point value such as $n-10^{-50}$ can therefore produce the wrong answer even when it appears extraordinarily accurate.

For primitive necklaces, Theorem~\ref{thm:two-candidate} provides a particularly simple implementation: bracket the leading solution using exact comparisons of $q^n/n$ with $Y$, then perform one exact comparison of $I_q(n_0)$ with $Y$. This avoids trying to choose one fixed-radical interpolation before the integer index is known.

\subsection{Computing coefficients without symbolic explosion}

The finite formulas make arbitrary weights well defined, but a direct symbolic expansion of all compositions is not always the fastest implementation. The following operations suffice for a stable truncated-power-series implementation:
\begin{equation}\label{eq:seriesexpalgorithm}
 f(t)=\ee^{g(t)}\quad\Longrightarrow\quad
 f_0=\ee^{g_0},\qquad
 f_n=\frac1n\sum_{j=1}^n jg_jf_{n-j},
\end{equation}
and, if $f_0\ne0$,
\begin{equation}\label{eq:serieslogalgorithm}
 g(t)=\log f(t)\quad\Longrightarrow\quad
 g_n=\frac{nf_n-\sum_{j=1}^{n-1}jg_jf_{n-j}}{nf_0}.
\end{equation}
Both follow by comparing coefficients in $f'=g'f$. Substituting these into
\eqref{eq:dm-recurrence} computes the inverse through any chosen finite weight without differentiating a large expression. The code also checks the independent nonrecursive formula \eqref{eq:finite-dm} against this recurrence through weight eight.

For generalized Galois numbers, compute the finite products $\pp{j}$ once, form the composition coefficients $A_{r,M}$ by truncated convolution, multiply by the appropriate theta constant in \eqref{eq:h-coeff}, take a truncated logarithm, and apply the inverse-coefficient routine. Only $r$ theta constants are required for any number of sectors at fixed rank $r$.

\section{Reproducible numerical checks}
\label{sec:numerics}

The archive contains \texttt{verify.py} and its generated
\texttt{verification\_report.json}. The supplied run passed \textbf{1,090 checks} using 110 decimal digits for analytic calculations and exact Python integers or rational numbers for the enumeration and threshold tests. The checks include the small integer values of the group, Gaussian, and $q$-Catalan families; generalized Galois identities for $r=2,3$; a weighted Rogers--Szeg\H{o} identity; inverse coefficient substitutions through weight eight; the finite nonrecursive coefficient formula; endpoint inversion; fixed-radical logarithmic-phase inverses; and the two-candidate threshold rule at rational test targets.

For the inverse tables below, a known real value $x_*$ is chosen and
$Y=F(x_*)$ is computed from the specified interpolation. The reported error is
\begin{equation}\label{eq:tableerror}
 E_M=\left|s+\sum_{m=1}^M d_m(s)t_0^m-x_*\right|.
\end{equation}
The displayed values are rounded; the JSON report retains additional digits. The truncation index $M$ always refers to the chosen small variable, so a zero weight still counts toward $M$. In particular the symplectic row uses $t_0=Q^s$, not the reduced variable $Q^{2s}$.

\begin{table}[htbp]
\centering\small
\setlength{\tabcolsep}{4pt}
\begin{tabular}{@{}lrrrrr@{}}
\toprule
Family & $x_*$ & $E_0$ & $E_2$ & $E_4$ & $E_8$\\
\midrule
$\GL$ & $8.25$ & $2.874\times10^{-4}$ & $2.674\times10^{-10}$ & $6.375\times10^{-16}$ & $6.908\times10^{-27}$ \\
Central Gaussian & $8.25$ & $5.738\times10^{-4}$ & $1.25\times10^{-9}$ & $7.171\times10^{-15}$ & $4.475\times10^{-25}$ \\
$q$-Catalan & $8.25$ & $4.578\times10^{-4}$ & $1.063\times10^{-9}$ & $6.282\times10^{-15}$ & $3.912\times10^{-25}$ \\
Symplectic & $8.25$ & $1.526\times10^{-7}$ & $1.983\times10^{-13}$ & $4.006\times10^{-19}$ & $2.455\times10^{-30}$ \\
Unitary, even sheet & $8.25$ & $9.558\times10^{-5}$ & $1.484\times10^{-10}$ & $3.491\times10^{-16}$ & $2.874\times10^{-27}$ \\
Unitary, odd sheet & $8.25$ & $9.589\times10^{-5}$ & $1.489\times10^{-10}$ & $3.5\times10^{-16}$ & $2.881\times10^{-27}$ \\
Galois $r=2$, $\eps=0$ & $16.25$ & $1.512\times10^{-3}$ & $4.131\times10^{-9}$ & $2.824\times10^{-14}$ & $2.62\times10^{-24}$ \\
Galois $r=2$, $\eps=1$ & $17.25$ & $1.007\times10^{-3}$ & $1.307\times10^{-9}$ & $4.276\times10^{-15}$ & $9.105\times10^{-26}$ \\
Galois $r=3$, $\eps=0$ & $18.25$ & $6.633\times10^{-3}$ & $3.316\times10^{-7}$ & $1.318\times10^{-12}$ & $3.45\times10^{-19}$ \\
\bottomrule
\end{tabular}
\caption{Quadratic-phase examples at $q=2$. Absolute inverse errors after selected exponential weights.}
\end{table}

\begin{table}[htbp]
\centering\small
\setlength{\tabcolsep}{4pt}
\begin{tabular}{@{}lrrrrr@{}}
\toprule
Family & $x_*$ & $E_0$ & $E_2$ & $E_4$ & $E_8$\\
\midrule
Fixed rank $k=3$ & $12.25$ & $6.913\times10^{-4}$ & $3.333\times10^{-11}$ & $3.334\times10^{-18}$ & $5.338\times10^{-32}$ \\
$q$-factorial, base $1/2$ & $12.25$ & $2.962\times10^{-4}$ & $2.556\times10^{-11}$ & $5.128\times10^{-18}$ & $3.908\times10^{-31}$ \\
Prime-power index $p=2$ & $16.25$ & $5.682\times10^{-3}$ & $2.537\times10^{-9}$ & $2.81\times10^{-15}$ & $2.754\times10^{-27}$ \\
Fixed radical $R=6$ & $36.25$ & $5.334\times10^{-6}$ & $5.334\times10^{-6}$ & $1.212\times10^{-9}$ & $3.321\times10^{-17}$ \\
\bottomrule
\end{tabular}
\caption{Linear and logarithmic phases. The last two rows use the exact $W_{-1}$ center.}
\end{table}

For the finite-limit inverse at $q=2$, take $x_*=16.25$ and
\[
 \eps=\mathcal A_{1/2}(2^{-16.25})
      =1.28310884632239305154\ldots\times10^{-5}.
\]
Truncating the convergent series \eqref{eq:endpoint-all} gives:
\begin{center}
\begin{tabular}{@{}rrrrr@{}}
\toprule
$E_0$ & $E_1$ & $E_2$ & $E_4$ & $E_8$\\
\midrule
$3.085\times10^{-6}$ & $1.414\times10^{-12}$ & $1.075\times10^{-18}$ & $7.137\times10^{-30}$ & $2.592\times10^{-52}$ \\
\bottomrule
\end{tabular}
\end{center}

As a separate check of the singular transition, set $\tau=1$ and retain
terms through $h^3$ in \eqref{eq:central-double-asympt}, omitting its much
smaller modular factor. The observed errors are consistent with the proved
$O(h^5)$ remainder:
\begin{center}
\begin{tabular}{@{}rrr@{}}
\toprule
$h$ & Exact $\log H_{1,1}(1/h)$ & Absolute error through $h^3$\\
\midrule
$.2$ & $8.475418514412$ & $4.856\times10^{-7}$ \\
$.1$ & $17.97956130094$ & $1.551\times10^{-8}$ \\
$.05$ & $37.31564489494$ & $4.874\times10^{-10}$ \\
\bottomrule
\end{tabular}
\end{center}

The sparse symplectic and fixed-radical rows also verify the predicted missing
weights. Such tests are more informative than a single agreement of the
leading approximation: a wrong small parameter, a lost parity, or a missing
curvature term usually changes the next nonzero error by several orders of
magnitude.

\section{Connections and further directions}
\label{sec:conclusion}

\subsection{A unified coefficient calculus}

The examples exhibit the same underlying structure at three levels. Finite $q$-products yield explicit logarithmic coefficients by \eqref{eq:product-log-coeff}. Theta-weighted flag sums yield them by \eqref{eq:h-coeff} and \eqref{eq:Galois-b}. Fixed-radical divisor sums yield them from the logarithm of the finite polynomial $E_R$. In all three cases, inversion is a finite operation at every exponential weight, expressed by \eqref{eq:finite-dm} or \eqref{eq:general-derivative}.

The coefficients record two distinct nonlinear effects: differentiation of the dominant inverse phase, and multiplication or shifting of exponential sectors. The first effect introduces factors such as $D^{-1}$ and derivatives of $D$; the second produces the composition sums $C_{m,k}$. Keeping these mechanisms separate explains the universality of the formulas without suppressing family-specific cancellations.

\subsection{Exact discrete identities can improve analytic interpolation}

The Galois-number construction illustrates a useful general strategy. A finite sum is extended to a lattice sum because an infinite-product factor vanishes outside the original finite domain. After an Euler expansion, completing a quadratic form converts every coefficient into a theta constant on a translated lattice coset. The result is an entire coefficient-generating function and an explicit residue-sheet interpolation, not only a leading asymptotic estimate.

For a more general finite-field flag or partition-function problem, the analogous questions are concrete: does a product factor enforce the finite domain by zeros; does the dominant exponent have a positive definite quadratic form on the constraint lattice; and do the coefficient translations preserve finitely many lattice cosets? When all three conditions hold, the proof of Theorem~\ref{thm:theta-tail} suggests a direct route to all sectors and then to their inverses.

\subsection{Problems left open by this organization}

Several extensions merit separate treatment. The first is a uniform theory in which the flag length $r$ grows with $n$: the lattice dimension and the exponential spacing $h/r$ then both vary, so the fixed-$r$ coefficient bounds are not uniform. The second is a complete resurgent analysis of the finite-size tail in the double-scaling limit; the exact modular product isolates one contribution but does not replace that analysis. The third is efficient interval evaluation of high-dimensional theta constants, allowing the explicit inverse series to serve as certified threshold algorithms rather than only high-precision approximations.

These questions do not affect the established fixed-parameter results. For fixed $q>1$ and fixed rank or flag length, the product and theta inverses above are convergent, their coefficients are explicit to every order, and their integer thresholds are controlled by the stated rounding or exact-comparison rules.

\appendix
\section{A compact Wolfram Language implementation}
\label{app:wolfram}

The following functions use the documented \texttt{QPochhammer},
\texttt{QBinomial}, and \texttt{ProductLog} conventions
\cite{wolframpoch,wolframqbin,wolframW}. The list \texttt{bs} contains
$\{b_1,\ldots,b_M\}$, without a leading zero.

\begin{lstlisting}
ClearAll[alpha, pInf, aTail, inverseCoefficient];
alpha[m_Integer, Q_] := Q^m/(m (1 - Q^m));
pInf[Q_] := QPochhammer[Q, Q];
aTail[Q_, t_] := -Log[QPochhammer[Q t, Q]];

inverseCoefficient[m_Integer, a_, lam_, D_, bs_List] :=
 Module[{u, bpoly},
  bpoly = Sum[bs[[j]] u^j, {j, 1, m}];
  -Sum[
    Coefficient[Expand[bpoly^k], u, m]/k *
     Sum[Binomial[k + j - 1, j] If[j == 0, 1, a^j] *
       (m lam)^(k - 1 - j)/
       ((k - 1 - j)! D^(k + j)), {j, 0, k - 1}],
    {k, 1, m}]
 ];

(* General linear group interpolation: work with its logarithm. *)
ClearAll[glLog, glInverseTruncated];
glLog[x_, q_] := x^2 Log[q] + Log[pInf[1/q]] +
                 aTail[1/q, q^(-x)];
glInverseTruncated[logY_, q_, M_Integer] :=
 Module[{h = Log[q], Q = 1/q, s, D, bs},
  s = Sqrt[(logY - Log[pInf[Q]])/h];
  D = 2 h s;
  bs = Table[alpha[j, Q], {j, 1, M}];
  s + Sum[inverseCoefficient[m, h, h, D, bs] q^(-m s),
          {m, 1, M}]
 ];

(* Exact integer examples; q is a prime power for field counting. *)
ClearAll[galoisNumber, qCatalanBinomial, irreducibleCount];
galoisNumber[n_Integer, q_] := Sum[QBinomial[n, k, q], {k, 0, n}];
qCatalanBinomial[n_Integer, q_] :=
  (q - 1) QBinomial[2 n, n, q]/(q^(n + 1) - 1);
irreducibleCount[n_Integer, q_Integer] /; n >= 1 && q >= 2 :=
  Total[(MoebiusMu[#] q^(n/#)) & /@ Divisors[n]]/n;

(* The leading inverse of q^x/x, on the large real branch. *)
ClearAll[necklaceCenter];
necklaceCenter[y_, q_] := -ProductLog[-1, -Log[q]/y]/Log[q];
\end{lstlisting}

Sequence indices are nonnegative; \texttt{irreducibleCount} requires a positive index. Inversion targets are positive and sufficiently large, with $q>1$ except where a base $Q\in(0,1)$ is explicitly used. The binomial-version $q$-Catalan function at $q=1$ is obtained by a limit. For symbolic calculations, take $q$ exact, such as $q=2$, and request numerical precision only after forming the expression. Near an integer threshold, verify candidate indices by exact comparison rather than relying on a machine-precision ceiling of \texttt{ProductLog}.

\section{Reproduction and a formula index}
\label{app:reproduce}

Compile with XeLaTeX until references stabilize (normally three runs):
\begin{lstlisting}
xelatex -interaction=nonstopmode combinatorial_transseries_extension.tex
xelatex -interaction=nonstopmode combinatorial_transseries_extension.tex
xelatex -interaction=nonstopmode combinatorial_transseries_extension.tex
\end{lstlisting}
The source uses Computer Modern Unicode fonts when available, with Latin Modern fallbacks. No font files are included. To reproduce the numerical report, install \texttt{mpmath} in a Python environment and run:
\begin{lstlisting}
python verify.py --output verification_report.json
\end{lstlisting}
No internet access is used by the verification program.

The finite nonrecursive inverse coefficient is \eqref{eq:finite-dm}; the general analytic-phase version is \eqref{eq:general-derivative}. Product families are specified by \eqref{eq:product-log-coeff}. All generalized Galois theta-tail coefficients are given by \eqref{eq:h-coeff}, their logarithmic coefficients by \eqref{eq:Galois-b}, and their full inverse sheets by \eqref{eq:Galois-inverse-full}. The endpoint inverse is \eqref{eq:endpoint-all}. Arithmetic fixed-radical sheets are \eqref{eq:fixedradical-function} and \eqref{eq:radical-inverse}; the uniform integer threshold rule is \eqref{eq:two-candidate}. The simultaneous $q\to1$ phase is \eqref{eq:Sphase}, with inverse coefficients \eqref{eq:crossover-all}.

\begin{thebibliography}{99}
\small
\bibitem{base}
Vladimir Reshetnikov, \emph{Transseries and Inversion}, consolidated article in the
\texttt{ProveIt} repository, directory
\href{https://github.com/VladimirReshetnikov/ProveIt/tree/main/Analysis/FabiusFunction/docs/semi-formalized-research-frontiers/drafts/series-and-transseries/Transseries_And_Inversion}{\texttt{Transseries\_And\_Inversion}}.
Source inspected 4 September 2026; main TeX blob identifier
\texttt{aec3c70d4793e286b462e7d8894dede7c754995e}.

\bibitem{dlmfq}
NIST Digital Library of Mathematical Functions,
\href{https://dlmf.nist.gov/17.2}{\S17.2, Calculus: $q$-Pochhammer symbols, Gaussian coefficients, and the $q$-binomial theorem}.
In particular equations 17.2.1, 17.2.3, 17.2.27, and 17.2.35--17.2.38.

\bibitem{dlmfqmod}
NIST Digital Library of Mathematical Functions,
\href{https://dlmf.nist.gov/17.2.E6_1}{equation 17.2.6\_1, modular transformation of the Euler product};
see also \href{https://dlmf.nist.gov/23.18}{\S23.18, Modular Transformations}.

\bibitem{dlmfW}
NIST Digital Library of Mathematical Functions,
\href{https://dlmf.nist.gov/4.13}{\S4.13, Lambert $W$-Function}.

\bibitem{wolframqbin}
Wolfram Research, \href{https://reference.wolfram.com/language/ref/QBinomial.html}{\texttt{QBinomial}}, Wolfram Language documentation.

\bibitem{wolframpoch}
Wolfram Research, \href{https://reference.wolfram.com/language/ref/QPochhammer.html}{\texttt{QPochhammer}}, Wolfram Language documentation.

\bibitem{wolframW}
Wolfram Research, \href{https://reference.wolfram.com/language/ref/ProductLog.html}{\texttt{ProductLog}}, Wolfram Language documentation.

\bibitem{oeisGL}
OEIS Foundation, \href{https://oeis.org/A002884}{A002884}, order of the general linear group $\GL(n,2)$.

\bibitem{oeisflags}
OEIS Foundation, \href{https://oeis.org/A005329}{A005329}, products $\prod_{j=1}^n(2^j-1)$.

\bibitem{oeiscentral}
OEIS Foundation, \href{https://oeis.org/A006098}{A006098}, central Gaussian binomial coefficients at $q=2$.

\bibitem{oeiscatalan}
OEIS Foundation, \href{https://oeis.org/A015030}{A015030}, binomial-version $q$-Catalan numbers at $q=2$.
The asymptotic entry by Vladimir Reshetnikov is dated 26 September 2021.

\bibitem{oeisSp}
OEIS Foundation, \href{https://oeis.org/A003923}{A003923}, orders of $\Sp(2n,2)$.

\bibitem{oeisU}
OEIS Foundation, \href{https://oeis.org/A070953}{A070953}, orders of the general unitary groups over $\F_4$ in the convention $\mathrm{GU}(n,2)$.

\bibitem{oeisGalois}
OEIS Foundation, \href{https://oeis.org/A006116}{A006116}, numbers of subspaces of $\F_2^n$.

\bibitem{kousidis}
Stavros Kousidis, \emph{Asymptotics of generalized Galois numbers via affine Kac--Moody algebras}, 2011,
\href{https://arxiv.org/abs/1109.2546}{arXiv:1109.2546}.

\bibitem{bliem}
Thomas Bliem and Stavros Kousidis, \emph{The number of flags in finite vector spaces: Asymptotic normality and Mahonian statistics}, 2011,
\href{https://arxiv.org/abs/1109.4624}{arXiv:1109.4624}.

\bibitem{oeisI}
OEIS Foundation, \href{https://oeis.org/A001037}{A001037}, degree-$n$ irreducible polynomials over $\F_2$, primitive binary necklaces, and binary Lyndon words.

\bibitem{oeisnecklaces}
OEIS Foundation, \href{https://oeis.org/A000031}{A000031}, binary necklaces without reflection equivalence.

\end{thebibliography}
\end{document}
